\begin{savequote}[450pt] %指定引言宽度
    \fontsize{8pt}{8pt} %两个参数分别指定字号和行间距
    明明有星空和浪潮，看起来却像在触礁。\footnote{\bilibili{BV1r4411A77J}；\netease{1865377022}。}
    \qauthor{闹闹丶\&果汁凉菜《无题》}
    曾幻想长大展翅飞翔，如今一副瘦肩膀。\footnote{\bilibili{av10154377}。}
    \qauthor{vsinger团队\&和田野《未来的我》}
    略过谎言的微妙，陷入圈套。所有该被保护的人竟然全都被忘掉。\footnote{\bilibili{av9100200}。}
    \qauthor{时之歌project\&冥凰\&苍白《迷局》}
\end{savequote}

\chapter{能力与成长\label{chap:成长}}

\addtocounter{section}{-1}
\section{基础讨论与用语}
本章内容涉及对\note{人的意识演化基本模型}的更深入讨论。本章会使用到一些之前篇幅中介绍过的概念，但涉及的重点和惯用措辞会有所区别。为了方便读者查阅和使用，现列举如下：
\begin{itemize}
\item \indicate{行为模式}：我们将\indicate{一个人在某些条件下，所有可能被触发的行为链}称为这个人在这些条件下的\indicate{行为模式}，具体讨论见\hyperref[sec:行为模式]{2.2.2小节}。在本章中，我们倾向于将\indicate{这些行为链}本身直接称为行为模式，而不硬性要求其由某个外部条件触发。
\begin{explain}
此处的侧重点改变不会影响“行为模式由某些\indicate{条件}决定”和“行为模式包含了\indicate{所有}可能触发的行为链”这两点定义。前者成立的原因是“我们可以将这些行为链本身视作触发行为模式的条件”，而后者成立的原因是“行为链会触发的东西，原来的条件也会（因为触发了行为链而）触发”。
\end{explain}
\item \indicate{竞争}：我们将\indicate{有多个行为}\footnote{或行为链、行为模式。}\indicate{同时触发，但其中只有部分行为被执行}的现象称为这些行为间的\indicate{竞争}，具体讨论见\hyperref[sec:并行与竞争]{3.4.1小节}。在本章中，我们倾向于只考虑\indicate{一个刺激触发的多个行为}之间的竞争，并将其称为由该刺激\indicate{引起}/\indicate{生成}/\indicate{形成}/\indicate{导致}的竞争。
\begin{explain}
此处的侧重点改变不会影响前文中对竞争机制的分析。相关的概念，如\indicate{覆盖}与\indicate{外溢}依然有效。“单一刺激形成的竞争”和“（某一时间点处）所有被触发行为之间的竞争”相比，最大区别在于前者只包含后者的部分行为。但在大多数情况下，在前者中被覆盖的行为，在后者中也会被覆盖。对于将少数例外情况，我们可以通过“将其视为刺激不同的另一个行为”等方式来规避。
\end{explain}
\end{itemize}
以上的两种处理可以方便我们讨论行为模式之间的覆盖关系。%，具体展开见\hyperref[sec:不可区分性]{5.1节}。

\labelsection{不可区分性}

\subsection{视角与不可区分性}
我们将\indicate{一类现象}称为一个\source{视角}/\indicate{关注点}/\indicate{可见范围}，将“一个环境/场景/事件中，属于某一个视角的所有现象”称为“这一视角在这一环境/场景/事件中的\indicate{可见}的部分/现象”，或是“这一环境/场景/事件在这一视角下的\indicate{可见部分}/\indicate{可见现象}”。
\begin{explain}
在不引起歧义的情况下，我们也将这一类现象直接称为一个\indicate{视角}。

视角可以有很多不同的类型，包括但不限于：
\begin{itemize}
\item 某人/某渠道在当前已知的信息：这可能包括“自己能回忆起来的东西”、“迄今为止记录并整理归档的数据”、“另一个人心理活动以外的行为”等。在不引起歧义的情况下，我们也将这个人（这个渠道）直接视为视角，并采取“在某人的视角下的信息”、“在外界看来”一类的措辞。
\item 会触发另一类现象的事件：这可能包括“某种学科的讨论范围（如化学只考虑化学变化）”、“某个目标的分解步骤”、“某种（可能外溢的）条件限制（如‘虽然闯红灯不一定出车祸，但不要闯红灯’）”等。我们会用“在经济/心理/化学方面”、“对某个目标来说”、“为了避免某个后果”等措辞来指代这些内容。

在一些时候，我们可以不区分“可见现象”和“可见现象的后续影响”，如“只要粉尘浓度合适，点火就会爆炸，粉尘是刻意扬起的还是不小心泄露的不重要”、“超市收银员不关心钱是哪里来的，只需要给钱就能买到东西”、“不管是追逐打闹、纯粹无知、赶时间还是被推倒，只要上了马路就更可能出车祸”等。
\item 会触发某人特定行为的事件：此类视角包含于上一类。在实际使用中，否定性叙述更加常见：某个现象无法触发A的行为（即“A的视角中不包含此现象”），从而事实上可以区分的内容对A来说不可区分。对此类视角的研究是本章的重点。

其应用场景包括但不限于：A因为认识不足（如语言不通），而无法识别某些差异（B所有话看起来都是叽里咕噜的）；A因为具有某些认知（如“B这人性格就这样”），而将不同的行为视为没有差别（“忍忍就过去了”）；A因为具有某些外溢行为（如“看见蛇就害怕逃跑”），而无法完成某些简单目标（“去井里打水”）......
\end{itemize}
在不同视角下，同一环境/场景/事件的可见部分不一定相同。比如，一个人可以从环境中发现更多现象，但是这个人身上的某个行为只会因其中一个现象而触发。人身上会有多个有意义\footnote{此处的“有意义”，指“在研究相关问题时可以起到作用”。本指南中的“相关问题”指“人身上的意识现象”。}的视角。因此，本指南此处不将“视角”定义为“一个人的认知”（甚至不是“一个行为模式的认知”），而定义为“一类现象”。
\end{explain}
在某一视角下，我们将“在某一环境中，有多个不同的过程/原因，都可以导致相同的可见现象”的现象称为“这些过程/原因在这一视角/这些可见现象下\indicate{不可区分}/\indicate{没有可见的差别/区别/不同}”，或是“这些过程/原因对这一视角/这些可见现象\indicate{没有区别}/\indicate{无关紧要}”，或是“这一视角/这些可见现象\indicate{不关心}/\indicate{不管}这些过程/原因，\indicate{无法发现}这些过程/原因的不同”，或是“其中一种原因在这一视角下与另一种原因\source{等价}/\indicate{等效}”，或是或是“其中一种过程/原因在这一视角下\indicate{等价于}/\indicate{等效于}另一种过程/原因”。
\begin{explain}
\trained 有可能更习惯“两种环境在可观测意义下等价”之类的表述，不过我们此处仍然选择比较贴近日常表达的措辞。

此处定义的不可区分性，只关注可见现象的成因，而不关注可见现象产生的影响。
\end{explain}
对于一些在视角A下不可区分，在视角B下可以区分的过程/原因：
\begin{adjustwidth}{0.05\linewidth}{0.05\linewidth}
\hspace*{2em}如果这些过程/原因会导致\indicate{不同}的结果，那么就将“视角B的可见现象中，不属于视角A的部分”称为“视角A关于这些结果的一种\indicate{有效信息}”，或者是“视角A\indicate{区分}这些结果的一种\source{必要信息}”。

如果这些过程/原因会导致\indicate{相同}的结果，那么就将“视角B的可见现象中，不属于视角A的部分”称为“视角A关于这些结果的\indicate{无效信息}”。
\end{adjustwidth}
\vspace{-1.0em}
\begin{framed}
    \vspace{-0.2em}
    \noindent 严格定义如下：

    对于两个现象集合（即环境/场景/事件）X和Y，我们将X和Y在视角A下的可见现象分别记为$\mathrm{X_A}$和$\mathrm{Y_A}$，在视角B下的可见现象分别记为$\mathrm{X_B}$和$\mathrm{Y_B}$。它们满足$\mathrm{X_A}=\mathrm{Y_A}$，$\mathrm{X_A}\subset \mathrm{X_B}$，$\mathrm{Y_A}\subset \mathrm{Y_B}$（此时，“视角B的可见现象中，不属于视角A的部分”就是$\mathrm{X_B}-\mathrm{X_A}$与$\mathrm{Y_B}-\mathrm{Y_A}$，其中符号“$-$”代表差集）。

    我们分别将X和Y的一个子集$\mathrm{X_O}$和$\mathrm{Y_O}$称作X和Y各自的“原因/过程”。这两个子集满足$\mathrm{X_A}$和$\mathrm{X_B}$由导致，$\mathrm{Y_A}$和$\mathrm{Y_B}$由$\mathrm{Y_O}$导致（此处默认因果性存在）。

    我们分别将$\mathrm{X_B}$和$\mathrm{Y_B}$的一个子集$\mathrm{X_C}$和$\mathrm{Y_C}$称作X和Y各自的“结果”（它们自然也分别由$\mathrm{X_O}$和$\mathrm{Y_O}$导致）。如果$\mathrm{X_C}=\mathrm{Y_C}$，我们就称“这些过程/原因会导致相同的结果”、“$\mathrm{X_B}-\mathrm{X_A}$与$\mathrm{Y_B}-\mathrm{Y_A}$是无效信息”；如果$\mathrm{X_C}\ne \mathrm{Y_C}$，我们就称“这些过程/原因会导致不同的结果”、“$\mathrm{X_B}-\mathrm{X_A}$与$\mathrm{Y_B}-\mathrm{Y_A}$是有效信息”。

    对于多个现象集合，上述定义中“在视角A下不可区分”指“它们两两不可区分”，“在视角B下可以区分”指“它们两两可以区分”。
    \begin{explain}
    $\mathrm{X_C}$和$\mathrm{Y_C}$是任意取定的，它们的唯一要求是“对视角B可见”。

    $\mathrm{X_O}$和$\mathrm{Y_O}$不要求对A或B可见（从而，我们可以规避一些过于细节和复杂的本体论讨论）。%它们中的每个元素分别都可能存在或不存在于$\mathrm{X_B}$和$\mathrm{Y_B}$中；它们的每个公共元素都可能存在或不存在于$\mathrm{X_A}$和$\mathrm{Y_A}$中。
    \vspace{-0.2em}
    \end{explain}
\end{framed}
\begin{explain}
视角B的可见现象可能包含结果本身（比如视角B是“视角A又过了一段时间，观察到了结果”的情况）。我们不在定义中排除这种情况，但将结果本身称为结果的\indicate{平凡信息}。如果视角B去除平凡信息后仍然能区分结果，那么我们就这些能用于区分的现象称为结果的\indicate{非平凡信息}。视角A的可见现象也可能包含结果本身，但只能包括相同的结果。

定义中提到了“一种”有效信息，是因为可能存在不止一种有效信息。这些有效信息之间可能没有公共部分，而即使有公共部分，两个有效信息的交集也不一定仍然是有效信息。从而，“视角A区分结果的必要信息”不具有唯一性，不是一个良定义的概念。无效信息的对应定义中没有“一种”，是因为无效信息不需要用于区分。不过，在不引起歧义的情况下，本指南也会略去“一种”二字。
\end{explain}

\subsection{不可区分的环境}
在某个环境中，我们将\indicate{一个行为模式触发的所有行为中，属于某个视角的部分}称为“这个行为模式在这个视角下\indicate{可见}的部分”，或是“\indicate{在这个视角下可见}的行为”，或是“在这个视角下的\indicate{可见行为}”。

如果这些可见行为本身构成一个行为模式，那么就将其称为（在这个视角下的）\indicate{可见行为模式}。
\begin{explain}
此处定义的可见，比4.4.2小节中定义的\note{可见}更加普遍，4.4.2小节中的定义，是将此处定义中的视角取为“行为模式A可以稳定解读的行为”时的特例。

依照上一小节的定义，可见行为同样可以被称为视角。但在后文中，我们将会使用可见行为来称呼，以避免产生歧义。相对地，后文讨论行为模式时，使用到的“视角”一词总是指代“某一类行为”。

同一个行为模式在不同的视角下，可见的部分会有所不同。一些常见可能性在上一小节已有列举讨论，此处不再赘述。

“可见行为本身构成一个行为模式”的意思是“这些行为已经包含了所有可能触发的行为链”，即“只要行为A是可见的，并且它（或者它造成的影响）可以触发行为B，那么行为B就也是可见的”。

关心可见行为模式，可以不遗漏地研究某个原因的所有结果。这对后续的理论建构极为重要。
\end{explain}
我们将\indicate{在某个视角下，一个行为模式在两个（或多个）环境中，可见的部分相同}的现象称为“在此视角下，这两个（多个）环境对该行为模式\indicate{不可区分}”。
\begin{explain}
根据视角选择的不同，环境之间的不可区分性也会不同。如果我们选择了极端常见的行为（如“呼吸”、“心跳”等，并且不计“屏住呼吸”“心跳加速”等细节）作为可见行为，那么任意两个环境都是不可区分的。

在我们后续的理论建构中，我们经常将可见行为选为“某一特定的行为模式”（此时它是一个可见行为模式）。这主要是为了描述如下现象：在接触不同环境时，一开始触发的行为链是不同的，但在经历了一些行为后，行为链会趋同。为了方便讨论，我们暂且将行为链趋同前的部分称为“前置阶段”。

在前置阶段，不同的环境会触发不同的行为。这些行为有可能是一些下意识反应，也有可能是一些零碎的想法，也有可能是一些有意的举动。这些行为会产生一些刺激。刺激有可能是行为本身，有可能是自己得出的结论，有可能是外界产生的变化。这些刺激又会触发一些新的行为。

一些刺激会触发我们所关心的可见行为模式。这些刺激有可能直接来自环境，也有可能如上一段所述，来自自身与环境的交互。如果关心刺激的具体来源，反而会不方便研究可见行为模式。

两个不同的环境如果能触发同一个可见行为模式，则一定涉及某种形式的覆盖。这有可能是“因为学到/总结出了知识，所以可以用统一的思路，分析两个表面看起来不同的现象”，也可能是“只是看到了某个很突出的特点，然后某个行为模式外溢到了别处”。覆盖有正面的，也有负面的，而我们可以不包含感情色彩地对这种现象展开统一的研究。并且，我们也应该不带感情色彩地研究，带有感情色彩反而有可能使我们得出有偏颇的结论，并且影响我们进一步的理解和判断，影响我们处理和解决问题的能力。

之前的篇幅中已经讨论过一些不可区分的环境。比如说，当我们将视角取为“获得某个认知”时，不同的资源是不可区分的。虽然它们会触发不同的理解途径，但这作为前置阶段，不在我们的可见范围之内。相反，如果我们将“其它相关的认知”“掌握的深度和系统性”等现象一并纳入视角，那么不同的资源就可以区分。
\end{explain}

\subsection{不可区分的行为模式}
在某个环境中，我们将\indicate{一个人（行为模式）接触这个环境之后，产生的所有现象中，属于某个视角的部分}称为“这个人（行为模式）在这个视角下\indicate{可见}的后续/后续现象”，或者是“这个视角下的\indicate{可见后续}”。
\begin{explain}
定义中出现了“人”，这是因为我们总是可以注意到“一个人出现在了环境中”，但未必知道其具体触发了什么行为模式。这样定义可以简化后续的讨论。

任何现象都可以是可见后续，它与行为模式不直接相关。但当我们使用“可见后续”的称呼时，总是在研究行为模式。

可见后续的定义中仅涉及到了“先触发行为模式，后产生现象”的时间顺序，而没有涉及因果顺序。这是因为，我们经常不能确认行为模式和后续现象是否具有因果性。在很多时候，我们也缺少一个合适的基准，无法使用“当这个行为模式不在场的时候，会发生什么”之类的思路。
\end{explain}
我们将\indicate{在某个视角下，两个（多个）行为模式接触了同一个环境后，可见的后续相同}的现象称为“在此视角下，这两个（多个）行为模式对该后续\indicate{不可区分}”。
\begin{explain}
我们不止可以在多个行为模式之间判断是否可区分，同时也可以在多个任意类型的成因之间判断是否可区分。但本指南主要研究人的意识现象，此处也仅针对“理解行为模式”而强调定义不可区分性。

根据视角选择的不同，行为模式之间的不可区分性也会不同。如果我们选择了极端常见的现象（如“地球在转动”）作为可见后续，那么任意两个行为模式都是不可区分的。

之前的篇幅中已经讨论过一些不可区分的行为模式，比如\hyperref[sec:情绪、感受与其分析与处理方法]{3.7节}和\hyperref[sec:指挥现象的一个实例解读]{4.5.4节}。读者可以以此为例，对照理解此处的概念。

如果我们采取“所有可能的行为链”来定义行为模式，则可以直接考察“不同的行为链可以产生同样的后续结果”的现象；但如果我们仅将行为模式视为“在某一个环境下，产生了某种后续现象”，那么此处定义的不可区分就没有实际意义，当我们将后续现象选为“这个人的一些举动”的时候尤其如此。

这是因为，我们只选取了一个行为模式的可见部分，而没有对行为模式本身展开研究。此时，我们需要使用“一个人在环境P和环境Q中都会触发行为模式A（并且P和Q对A不可区分），另一个人在环境Q和环境R中都会触发行为模式B（并且Q和R对B不可区分），A和B在环境Q中不可区分”的视角来看待问题。

虽然在定义行为模式的不可区分性时，我们在为后续现象溯因，但因为我们实际关心的是行为模式，这一概念也将起到一些预测的功能。
\end{explain}
以下是一些常见的“行为模式不可区分”的情况：
\begin{explain}
以下的例子中，我们将人称为A。在有必要的情况下，将视角称为B。
\begin{itemize}
\item \indicate{不存在与未触发}：当可见后续为“A没有某个特定行为”时，既有可能是“A不知道要做什么”（即“行为不存在”），也有可能是“A没有识别到信号”（即“行为未触发”）。此时有三种可能的情况：“知道要做什么但不知道现在要做”、“知道现在要做但不知道要做什么”、“既不知道做什么也不知道现在要做”\footnote{如果此时不存在信号，那么仅有“A没有对应的行为”和“A有对应的行为但是未触发”两种情况。}。

此类现象的一种常见现实情境是“B在教A，发现A不会做题”。此时，“A没有掌握知识”和“A不会读题”是两种不同的情况，需要两种不同的讲法。讲法如果错配，就无法教会A。

\item \label{para:被覆盖与缺失}\indicate{被覆盖与缺失}：当可见后续为“A没有某个特定行为”，并且我们可以确定“A接收到了信号”时，既有可能是“A不知道要做什么”（即“行为不存在”），也有可能是“A更想做其它的事情”（即“触发了行为，但其与其它行为产生了竞争，并且其它行为胜出”）。这两种情况是互斥的，但前者可以触发后者。

此类现象的一种常见现实情境是“B在教A，发现A没有在听”。此时，“A听不懂”和“A走神了”是两种不同的情况，前者需要补充更基础的知识，而后者需要提醒A专注，有必要时还需要暂停，先行处理A遇到的问题。方法如果错配，就不会起到效果。

\item \indicate{动机与影响}：当可见后续为“A某个行为的某一结果（影响）”\footnote{我们不需要假设“视角可以确定行为和影响之间存在因果关系”。视角可以因为一些外溢认知，错误地认为两个不相关的东西之间存在因果关系。这不影响我们后续的分析和判断，我们仍然可以得出“此时无法判断A的动机”的结论。}时，A既有可能“为了这一结果从而行动”，也可能“因为别的刺激而行动”。“别的刺激”有可能是此行为的另一个影响，可能是此行为实际上没有产生（但A因为希望这一影响出现而行动）的影响。A也有可能完全没有考虑影响，而是触发了某个不可控行为（比如“转述别人的行为”）。

当我们使用“A的内心深处其实渴望着xxx”一类的描述时，经常会混淆动机（目的）与影响（结果）：很多情况下，我们能观察到的现象是“A做了某些举动，B做了某些举动后，A感到开心/满足/停止生气/......”。此时，B的这些举动有可能满足了A的需求，也有可能只是一次无关的覆盖，不能视为A的渴望。具体的讨论参见\hyperref[sec:局部有效理论]{5.2.3小节}。

\item \indicate{是否自知/理解/可控}：当可见后续为“A的某个行动”或是“A某个行动的某个影响”时，A既有可能知道“自己做了这件事”，也可能不知道；既可能知道“自己为什么会做这件事”“事情为什么会变成这样”（也即“理解行动/行动的影响”），也可能不知道；既可能是出于自己决定而做，也可能不是。

这些情况直接决定了“A能否根据某些目标，调整自己的行为”。目标有可能来自自身要求，有可能来自外界要求，但它们不一定能够生效。人不一定因为“我应该做到xxx”“你必须做到xxx”就一定能做到对应的事情，其中可能有不可忽视的能力缺失。具体的讨论参见\hyperref[sec:假性成长]{5.6节}。
\end{itemize}
\noindent 除此之外，我们还会在一些比较技术性的地方使用不可区分的概念，比如说“对于某个结果来说，‘当前的做法’和‘什么都不做’不可区分”（以此来说明“这是某个目标的无效操作”等事情）等。
\end{explain}

\labelsubsection{广义行为与多重触发}
我们将\indicate{除了第一个行为，每一个行为都由之前的行为触发；除了最后一个行为，每个行为都触发了后续行为}的行为链\footnote{对于\rigorous，此处的行为链必须只由有限个行为组成。更严格的定义是：我们可以将行为按触发关系画成一张有向无环图，这张图有唯一的源点（只有出度）“第一个行为”和唯一的汇点（只有入度）“最后一个行为”，并且对于所有其它点，都存在一条从第一个行为到最后一个行为的路，经过这个点。}称为一个\indicate{复合行为}或是\indicate{可分行为}，将“第一个行为的刺激”的视为触发这个复合行为的刺激/信号，将“最后一个行为”视为这个复合行为执行的行动。我们将其中的每一个行为称作这个复合行为的一个\indicate{步骤}。

相对地，我们将\indicate{一个行为}称为\indicate{原子行为}或是\indicate{不可分行为}。我们将\indicate{原子行为和复合行为}统称为\indicate{广义行为}。在不引起歧义的情况下，我们也将广义行为称作\indicate{行为}。
\begin{explain}
此处的“触发的刺激”和“执行的行动”与\hyperref[sec:并行与竞争]{3.4.1小节}中的含义相同，都是将行为拆分成了“原因”和“结果”两方面。

\note{行为链}的定义是“（在时间上）连续触发的行为”，不要求行为之间有因果关系。此处对复合行为的定义主要是为了区分“因为某个持续存在的刺激，而重复地触发某种行为，但这些行为之间没有因果关系”的情况。复合行为中不一定是链式的“前一个触发后一个”，也可以有分叉，但整体结构是“单一原因（第一个行为的刺激）导致了单一结果（最后一个行为）”。在要求不是那么严格的情况下，我们也可以允许将“某个行为导致了外界的某个结果，而这一结果又产生了刺激，触发了新行为”纳入复合行为中，这种处理会方便我们的讨论。

此处的“可分行为”和“不可分行为”指的是“能否分成更小的组成部分”。由于它和“不可区分性”用字相似，容易产生歧义，本指南将不使用这一对词汇，而是使用更清晰的“复合行为”和“原子行为”。“原子”是计算机领域的借词\footnote{对于\rigorous，计算机领域中的“原子操作”“原子指令”等概念，也是来源于哲学领域的借词，表示“（物质的）不可分单位”。本指南不做深入的词源辨析。考虑到“原子”一词的“不可分单位”概念可以互通，并且计算机领域的概念确实更为接近，故使用“计算机领域的借词”的表述。此处说是哲学方面的借词，而不是化学/物理方面的，是因为原子的“不可分单位”义项先于“由质子、中子、电子组成的物质实体”义项而存在，后者实际上是一个错误的命名（原子并没有“不可分”的性质），只不过用久了就成习惯了。}。
\end{explain}
\indicate{当视角中仅能观察到“信号和行动先后出现”时，复合行为和原子行为不可区分；不同的复合行为不可区分。}
\begin{explain}
当我们将“信号”视为“环境”，“行动”视为“可见后续”时，这里的叙述就是“行为模式的不可区分性”。

该结论是此处引入这一对概念的主要原因。如果我们只关心行为，其实不用给它另外起个名字。此处的视角可以任意取定，具体可以参考\hyperref[def:视角]{前文}的讨论，此处不再赘述。

除了复合行为和原子行为，信号和行动也有可能没有因果关系，只是在时间顺序上先后出现。这种情况在这种视角下也不可区分，但我们此处暂时不关注这种情况。

复合行为必然会产生某些除了信号和行动以外的特征（比如说思路），但是这些特征不一定是可见的；而在一个行为触发时，即使观察到了一些额外的特征，也不一定说明这一行为就是复合行为，这一特征有可能是这一行为的无关现实（比如由同一个信号产生的另一个行为）。这说明，在仅观察到“信号导致行动”的时候，我们既无法得到“这是一个原子行为”的准确判据，也无法得到“这是一个复合行为”的准确判据。

另一方面，我们可以从任何一个视角出发，定义该视角下的原子行为和复合行为\footnote{一些读者可能会将此处的论述视为“不可分性是相对的”，进而有“更深入，能分得更细的视角更优”的思路（比如，将“原子由质子、中子、电子形成”看成是“认识逐渐深入”的一个证据）。但是，“获得认识的时间顺序”、“认识的正确性和本质程度”、“认识的对现象所拆分的层次”是三件不同的事情，不同层次的内容需要用对应层次的视角研究（比如，“原子实际上可分”不影响“原子是最适合用于描述简单化学现象的理论工具”）。对系统论的相关讨论参见\hyperref[sec:复杂系统]{3.2节}，此处不再赘述。}。这种观点下，\indicate{“行为是否可分”是一个和视角相关的性质}。本指南主要研究意识现象，不深入生理过程，不将行为拆到“神经活动”层次。这可以有效避免一些会阻碍我们关心意识现象的，过度细节的讨论，比如说“抬起手总会先抬起一半，然后再抬起四分之一......”之类的阿基里斯式思路。

本指南仅关注可以产生\note{关联}的现象，将所有的\indicate{反射}视为原子行为。
\end{explain}
如果两个广义行为受到的刺激和执行的行动都相同，我们就称这两个广义行为\indicate{等价}或\indicate{等效}，或是将这些广义行为称作互相的\indicate{等价行为}/\indicate{等效行为}。
\begin{explain}
本指南在使用讨论广义行为时，默认采用“只关注受到的刺激和执行的行动”的视角。在此视角下，此处的等价/等效与\hyperref[def:等价]{前文}定义相同。

等价的广义行为是广泛存在的。常见的等价行为包括但不限于：
\begin{itemize}
\item 对同一件事有同样的称呼：不同人对同一个词汇的接触途径不同。不同的人虽然会将同样的事物用同一个词称呼，但有可能是根据完全不同的特征产生的判断。这种情况下不同的人如果展开进一步的沟通，就容易发现双方聊不到一块去，总是鸡同鸭讲。
\item 从同一个现象中获得认知：由于视角不同，不同的人可以从同一个现象中看到多方面特征，并且根据这些特征，独立地形成认识和行为。在这种情况下，“遇到某个现象，就采取对应的行动”，和所有的“遇到某个现象，发现对应的特征；遇到某个特征，采取对应的行动”都是等价行为。
\item 同一个行为链的不同简化：行为链也是现象，其中的每个行动都可以视为一个特征。我们可以以任意一个行动为刺激，任意一个行动为行动，产生一个新的行为。刺激在行为链中的顺序不一定要比行动更靠前，有可能产生因果性破坏。只要新产生的行为仍然能形成一个行为链，这就是原先行为链的等价行为。
\end{itemize}
\end{explain}
我们将\indicate{一个刺激触发多个等价的广义行为}的现象称为这一（广义）行为的\indicate{多重触发}，将其中的每一种广义行为称作这一多重触发的一个\indicate{触发途径}，将触发途径的个数\footnote{对于\rigorous，此处的准确表述为“不相互包含的触发途径的个数的最大值”。}称为这一多重触发的\indicate{数量}或者\indicate{重数}。
\begin{explain}
在不引起歧义的情况下，我们也会采用“n重触发”的表达，其中n是这一多重触发的重数。

% 多重触发在语义上也可以被称作“复合的多个行为”，但本指南中的“复合”已经用于指代行为链，不会用于多重触发。

多重触发现象是广泛存在的。我们可以从同一个现象中获取多个等价行为，这些行为之间不一定会产生竞争和覆盖，有时可以并行。一般来说，同一个行为模式触发次数越多，其简化行为也就会越多，产生的多重触发（无论是单一行为的重数，还是会被多重触发的行为数）也会越多。这在外部看来，是“可以熟练应用这一行为模式”。

除此之外，也可以巧合地从不同的来源中凑出等价行为。这种现象（在意识到以后）一般会带来强烈的感受和认知（比如“极大的震撼”、“重要的人生经验”、“灵魂的归宿”等），比如“殊途同归”“大道至简”等。

多重触发现象使得在一些情况下，即使我们做严格的实验，也无法区分原子行为和复合行为。比如说，如果信号A可以同时触发行为B和行为C，而行为B的某个影响又会触发行为C，那么当我们单独提取出行为B的这个影响时，C会触发，但“刺激A触发行为C是一个复合行为”的结论是错的。准确的结论需要专业的因果性分析，此处不展开论述。
\end{explain}

\section{竞争位置}
\labelsubsection{向前竞争与向后竞争}
我们将\indicate{在行为A的某次触发中，“行为A”或“行为A的某个影响”触发了行为B}的现象称为“行为A的此次触发中，行为B对行为A的\indicate{向后竞争}”。如果行为B覆盖了行为A，或是消除了行为A的影响，那么就将其称为“行为B对行为A的\indicate{向后覆盖}”。
\begin{explain}
这个定义看起来不像是人话。整个定义好像只是给“行为B被触发”另起了个名字，行为A只是个背景板，和竞争也没什么关系。

但其实这一定义的重点在于行为A。行为A的影响中，有一些是“我们需要有意处理”的，一般带有负面的价值判断。此时对应行为B的描述包括但不限于“收拾烂摊子”“擦屁股”等。定义中的“竞争”指“通过行为B而消除行为A的影响”，但触发B不保证能做到消除影响。这一定义主要用于对应“想要消除影响”的意图，而非“实际上消除了影响”的成果。

此处，为了避免引入过多新概念，同时为了与“向前竞争”形成更明确的对比，我们在以下两方面稍微滥用了一下术语：
\begin{itemize}
\item 竞争发生在“行为A的影响”和“行为B”之间，并不完全符合\hyperref[def:竞争]{前文}中“竞争一定会发生在行为之间”的定义。但在本指南的讨论范围之内，“行为A的影响”在绝大多数情况下，或者是行为A本身（此时它是一个持续时间很长的行为，如“心情低落”），或者是行为A触发的另一个行为C。在此时，“A/C和B之间的竞争”符合之前的定义，而“消除A的影响”也就变成了“用B打断行为A/C，实现覆盖”，符合覆盖的定义。

比如，行为A可以就是“想做行为C”的想法，而行为B是“决定是不是要做行为C”的决策。如果此处再滥用一下术语的话，可以得到“决策总是对动机向后竞争”的结论。

\item 行为A和行为B不一定属于同一个人，但这不影响我们关于“向后竞争”和“向后覆盖”的定义。如果X做了行为A，Y因此做了行为B，我们也将其称为“向后竞争”；如果行为B影响到了X使得行为A被打断，或是消除了行为A的影响，我们也将其称为“向后覆盖”。
\end{itemize}
\end{explain}
我们将\indicate{在行为A的某次触发中，“触发行为A的信号”或是“会产生这一信号的现象”同时触发了行为B}的现象称为“行为A的此次触发中，行为B对行为A的\indicate{向前竞争}”。如果行为B覆盖了行为A，或是消除了触发行为A的信号，那么就将其称为“行为B对行为A的\indicate{向前覆盖}”。
\begin{explain}
此处不使用“行为A与行为B产生了竞争”来定义“向前竞争”，是因为这无法与向后竞争区分，向后竞争中也有大量的“行为A与行为B产生了竞争”的现象。使用“行为B的触发条件”来区分二者，可以确保这它们不冲突。

对于我们上一段所举的“行为A是‘想做行为C’的想法，而行为B是‘决定要不要做行为C’的决策”的例子，行为B对行为A向后竞争，而对行为C向前竞争。

这个定义至少看起来确实是“行为A和行为B之间的竞争”，比上一个定义要更讲人话一点，但实际上，关注的重点依然在行为A上，也存在一定的术语滥用：
\begin{itemize}
\item 向前竞争时，也不一定会产生竞争，行为A和行为B可以并行。这个定义同样对应了“想要消除行为A”的意图，而不对应“实际上消除了行为A”的成果；

\item 行为A和行为B不一定属于同一个人，“Y的行为B发现信号以后，消除信号，使得X的行为A不触发”的情况也算是向前覆盖。
\end{itemize}
另一方面，如果行为A与行为B产生了竞争，有可能既不是向前竞争，也不是向后竞争，B可以由和A无关的另一因素触发。定义中的“向前”“向后”主要基于触发顺序（而不是完全时间顺序），即“B的触发条件出现时，是早于A触发时，还是晚于A触发时”。如果A和B不相关，那么自然不适用“向前/向后”的判断。
\end{explain}
如果行为B对行为A产生了竞争，但既不是向前竞争也不是向后竞争，我们就将其称为“行为B对行为A的\indicate{无顺序竞争}”。

我们将行为B对行为A“向前竞争”、“向后竞争”、“无顺序竞争”这三种情况，统一称为“行为B对行为A的\indicate{竞争位置}”。
\begin{explain}
更一般地，当我们试图打断某一个行为链时，“参与其中哪一个行为的竞争”也被称为“竞争位置”。不过，行为链的竞争位置可以由“对于行为链中的每一个行为，是向前还是向后”来确定，所以我们不独立地引入“行为链的竞争位置”的术语，仅在少数不引起歧义的情况下，为简化行文而使用。

“向前覆盖”与“向后覆盖”不是独立的竞争位置，它们的竞争位置分别是“向前竞争”和“向后竞争”。但当我们语境中需要关注“是否实现了覆盖”时，也会直接使用“向前覆盖”与“向后覆盖”作为竞争位置，并且将竞争位置也称为\indicate{覆盖位置}。

对于“行为A在一个行为模式中反复触发，而在某两次触发之间产生了行为B”一类的情况，需要根据“关注哪一次触发和竞争”，来判断行为B的竞争位置。比如，如果我们关注“脱离这一行为模式”一类的目标（比如说克服懒惰），关注“行为B打断了后续的行为A”，那么就算向前竞争；如果我们关注“纠正一个习惯”一类的目标（比如说戒酒），关注“还是有行为A”，那么就算向后竞争。

向前竞争与向后竞争是消除影响的两种策略。因此，本节中的事例分析会有一部分带有“发现问题，解决问题”的背景思路。

一部分读者在看到这里时，有可能会认为，这是治标与治本的区别：向后竞争只是对症下药，而向前竞争则从根源上解决问题。本指南不建议读者这么理解，原因有二：
\begin{itemize}
\item “治标”与“治本”本身是歧义很重的概念，我们难以得到“什么是标什么是本”的统一认识。虽然每个人都有模糊的感觉，但是模糊的感觉不足以让我们完成接下来的讨论，反而会增加更多的不解。

\item “治标”与“治本”具有价值与感情色彩方面的倾向，一般我们会认为治本更好。但向前竞争不保证比向后竞争更好。比如，对行为A的向前竞争在生效时，会不论正面负面地消除行为A的所有影响，而无法做到精确控制。比如“因为感觉自己会在社交中受创伤，就完全不交朋友”一类的处理方式。我们能很轻易地发现这是个过犹不及的处理方案，但这个模糊的感觉并不能带来有价值的建议和修正方法。
\end{itemize}
具体哪种情况更好，需要依照目标来确定。具体讨论见后文。
\end{explain}

\subsection{竞争位置的不可区分性}
\noindent \indicate{仅关注单一影响时，向前覆盖与向后覆盖不可区分。}\label{para:覆盖位置}
\begin{explain}
此处的“关注影响”指的是“在某种行为下会出现某种结果，在调整了行为后结果消失”的现象。“调整行为不影响结果”的情况不在此处讨论范围之内。

这一论述是真实情况的弱化版本。实际上，即使关注更多信息，向前覆盖与向后覆盖也经常不可区分。在后续的分析中，我们经常只关心“会作为触发某一行为的信号”的影响，而这大多数情况下是单一的，这一论述已经大体够用。故此处不使用更强的结论。

如果使用行为链来描述这条结论，那么会是“如果一个行为链会导致一个结果，那么无论打断行为链中的哪个行为，都有可能使后续行为不再存在，使最终的结果消失。不同的竞争位置不可区分。”这是因果关系的简单推论。本指南不引入行为链的竞争位置，此处仅用于辅助读者理解。
\end{explain}
存在\indicate{通过某次向后竞争，得到可以用于向前竞争的行为}的现象。
\begin{explain}
这一现象的准确描述是：先触发了一次行为A，并因此触发了某个行为模式C（C视作对A向后竞争），形成了行为B。后续再次触发A时，B可以参与对A的向前竞争。竞争位置是根据行为的触发顺序而定义的，而不是根据时间顺序而定义的，所以这不会违反定义。

如果我们适当扩充定义，那么这种现象可以视作一种\note{因果性破坏}。但由于论述过于复杂，我们后续行文中将不采用这一思路展开论述。

这一类现象极为常见，并且有多种表现形式，包括但不限于：
\begin{itemize}
\item \indicate{收集}：具有一些信息缺失，需要有意识地调研。比如“先试一遍附近的所有餐馆，看看哪家好吃”、“自己亲手做一遍流程/跟一遍新手教程”、“控制变量做实验以确认具体的相关性”、“观察别人的行为和影响”等。
\item \indicate{回忆}：没有需要刻意获取的信息，分析已有信息以获得结论。根据复杂程度的不同，可以有“刚一做就立刻发现自己做错了，下次就不做了”、“复盘自己在某方面的做法，从中发现问题并且修改”、“找到了自己的坏习惯，并且通过努力控制和练习培养了新的习惯”等。
\item \indicate{告知}：无法靠自身获取某些信息或是结论，需要依靠外部告知。比如“刚接触一个领域，不熟悉具体的评判标准，犯了熟手不会犯的错误”、“运动姿势不标准（容易引起损伤/会影响水平），需要另一套技术”、“对方持有和自己不同的观点，要求对方出示证据”等。
\end{itemize}
“从向后竞争中获得用于向前竞争的行为”可以认为是“学习能力”的一部分。会有将“从别人的严厉指责/辱骂中学习”当成“好学的模范”之类的例子。本指南不过度讨论这些以态度和情绪为主的例子。

需要注意的是，这并不是唯一的获取向前竞争的行为的方法。我们也可以从某些完全无关的事情中，巧合地得到向前竞争的行为，比如说“突然灵机一动，想通了一些事情”，或“在经历了另一些事情后，才明白了之前不知所云/嗤之以鼻的道理”；又或者，在经历了一些别的事情后，自身得到了潜移默化的改变，原先的行为不再存在，也就不需要考虑覆盖问题了。

但以上情况不可控，所以不会被人广泛采用，只会被归入“天赋”“阅历”等不可复制的特点之中。我们在后续的分析中，也仅会关注“从向后竞争中获取向前竞争的行为”的方式。这是唯一可控的解决问题的方式。
\end{explain}
在一个行为模式中参与向前竞争的行为，在另一个行为模式中可能参与向后竞争。
\begin{explain}
两方面的情况都有可能发生：
\begin{itemize}
\item \indicate{一个对X向后竞争的行为Y，在另一个行为模式中可能对X向前竞争}：参考上一段中的例子。%A观察到了对行为X的某个后果的负面评价Y（X和Y是A的行动还是别人的行动均可，Y需让A观测到，Y是否客观不重要），此时Y在对X向后竞争。如果A下次在做X之前想到了Y，那么Y就参与了向前竞争。更详细全面的例子
\item \indicate{一个对X向前竞争的行为Y，在另一个行为模式中可能对X向后竞争}：比如，A在做X之前，会有一个完整思路（其中包括某个判断Y）。A打算将思路教给B，但是B因为前置认知不足，不知道什么时候应该怎么办，无法掌握完整思路。B在做了X以后，会回忆起“A教了自己”的事情，并且回忆起一些判断（比如Y），然后后悔。
\end{itemize}
以上判断中的“向前/向后竞争”也可以换成“向前/向后覆盖”。如果我们同时关注竞争和覆盖，则会有更复杂的情况，比如：一个行为可能可以参与向前竞争，但无法向前覆盖，而只能向后覆盖。
\end{explain}
\indicate{仅关注一个行为A和一个可以与A向前竞争的行为B时，我们无法区分B是在向前竞争还是在向后竞争/无顺序竞争。}
\begin{explain}
这是本节第一个结论的简单推论。此时，我们将“出现了行为B”视作“关注的单一影响”。

这个结论看起来似乎自相矛盾，这主要是因为“可以与A向前竞争的行为B”的定义和指代不明。但我们此处无法给出一个更准确的定义。这个定义大体上指“实际上在某个人身上参与了向前竞争的行为”或是“属于某类‘可以向前竞争的行为’的行为”（比如“收益”、“态度”等），而实际上是指“任意类型的，会在想到‘如何更改行为A’时，同时想到的行为”，是个因人而异的主观定义（这不是在说“实际上大多数人根据‘向前竞争’这一特点找对应的行为，并且用于改正行为A”，因为在不知道“向前竞争”这一概念（以及等价的“治本”等概念）的情况下，无法做出以上操作。）。

当我们观察到“某次向后竞争以后，影响消失了”的现象的时候，我们很容易将其认为是“向后竞争本身有效”，但实际上也有可能是“向后竞争后，得到了可以用于向前竞争的行为”。此时，如果将同样的向后竞争应用在其它行为模式上，因为行为模式不同，不保证也能得到相同的行为。

在面对“可能可以参与向前竞争，但无法向前覆盖，而只能向后覆盖”的情况时，情况会尤为复杂。此时，即使我们考虑另外的影响，也只能得到“影响被消除”的信息，无法将这种情况与“向前竞争并向前覆盖”区分开来。这导致了大量的“知道了道理但还是做不好”等现象产生。
\end{explain}

\labelsubsection{局部有效理论}
\begin{explain}
本小节的注释过多。为了方便查阅，我们将注释放置在段尾而非页尾。
\end{explain}
如果一个认知中包括“某一现象具体到现象类型的归因”，我们就将这一认知（或者这一归因）称为关于此现象的一个\indicate{局部有效认知}/\indicate{局部有效理论}，将现象称为这一认知（理论）的\indicate{目标现象}。
\begin{explain}
定义中“具体到现象类型”指的是“给出了一个由现象类型组成的因果链\paranote{更清晰的表述是：我们可以将归因抽象为有向图（边的方向从原因指向结果），图的每个顶点是一个现象类型。用于分析目标现象时，就给每个现象类型代入相应的现象。}，并且该认知认为这是目标现象的唯一成因”。
\insertparanotes
“局部有效”这一名称的含义参考\hyperref[footnote:局部]{5.5.2小节的脚注14}。此外还有如下备选称呼\paranote{本指南不单独引入这些术语，仅在此列举，方便读者直观地理解这一概念。}：
\begin{itemize}
\item 短幅有效认知：指“已知信息（篇幅）仅有这一现象时，该成因是一种可能的情况”。

\item 向后有效认知：指“这一现象的后续影响与其成因不直接相关，可以基于此现象独立分析”。

\item 巧合有效认知：巧合指“原因分析错了”，有效指“能得到正确的结论（指这一现象的后续影响）”，和\hyperref[sec:巧合有效的沟通]{4.3节}中的“巧合”含义一致。

\item 形式有效认知：指“现实情况不属于理论的适用范围”，是一个从数学领域的借词。\paranote{对于\trained，我们可以注意到，当我们将某个（不收敛、扩充定义域、或其它类型的）形式运算过程视为归因，而运算结果视为现象的时候，形式运算就是运算结果的形式有效理论。此处的定义是严谨的。}
\end{itemize}
起“局部有效认知”的名字，不代表有“全局有效认知”一类的概念，也不意味着有“局部以外的情况”。此处的称呼仅描述客观特点，不包含负面的价值判断。事实上，如果我们已知的信息仅限“这一现象”本身，那么我们无从判断其是否还有可能有其它成因\paranote{本指南不深入“因果性是否存在”“归纳法是否有效”等认识论与逻辑学问题，不在此处过多展开讨论，仅稍作提及。}。
\insertparanotes
\end{explain}
我们以一个例子来说明这里的论述：\indicate{“别人这么对待你，其实是自己想要被这么对待”是一个局部有效理论。}
\begin{explain}
这个结论有很多种可能的来源：有可能来自各种心理学结论，也有可能来自社会阅历和人生感悟。它在社会中广泛传播，在各种心理学、成功学、（人际）关系学、文学的书籍、课程、论述、争辩中，以多种形式被不断提及，措辞可能非常不同，但是结论都一致。

我们在日常生活中，更多是在某些书籍/播客/自媒体听到，或自己总结出这个结论。因为篇幅、眼界、知识的限制，大多数情况下，它没有全套的背景分析，只是一个孤立结论。我们优先分析这种情况。
\end{explain}
以下举一种常见的情况作为例子：
\begin{explain}
当事人A\paranote{可能是书或者课程（在里面举了这个例子）的作者本人、心理咨询师的来访者，或其它身份。}每次在一个特定环境中\paranote{可能是在过年回老家，或是原生家庭，或是工作环境。}，都会遇到自己不想面对的人\paranote{可能是父母、亲戚、上司、同事、小孩等。这一般是无法躲开的人。}（记为B）。B会到处问\paranote{有可能只有一个很喜欢问的B，也有可能有好几个。}A各种各样的问题\paranote{以基础生活状况为主，比如学业、工作、婚姻、家庭等。}，A被迫招架，疲于应付。与此同时，B还会大肆评价A给出的回答，这些评价会让A非常反感\paranote{各种反感类型都有可能有，麻烦、无奈、反对、厌恶都很常见。}。
\insertparanotes
接下来会发生两件事情，一件是理论方面，一件是实践方面。这两件事没有明确的先后关系，不同的事例中会有不同的顺序。
\begin{itemize}
\item 实践部分：A\paranote{可能是巧合做出，或看到了别人这么做照着学，或通过分析得到应对方案等。}在B询问时，反过来询问B的相应情况\paranote{可能是直接问，甚至先发制人，或带有一些沟通技巧，如先回答/敷衍/转移话题再问。}。此时A会发现\paranote{或者吃惊地意外发现，或者毫不意外地验证了自己的分析。}，B开始详细地讲述自己的事情\paranote{可能是亲身经历，也可能是别人的故事；可能是二人的公共社交圈中的人，可能是B的个人社交圈。}，不再积极地询问。如果话题时间较长，还会有一些额外进展，比如“B讲自己想到/得到的看法/人生道理”\paranote{此处仅举一个例子，还可能有其它情况。}，而A也可能因此而得到B的良好评价\paranote{主要内容是“懂事”“合拍”等。不同的话题走向评价也会不同，但每个话题下的评价都是良性的。}。
\insertparanotes
\item 理论部分：A\paranote{可能是听别人说/从书里/网上看到，也可能是自己从实践经验中悟出来。}得到了“B问这个方面，是因为B想跟人说这个方面”的认识，并且基于此来理解与对待整件事\paranote{根据A原先对这一事件的态度不同，此处可能会涉及一些“克服情绪”或者“心理疏导”的步骤。我们这里所关注的现象是“A可以使用这一理论，而不会被自己的其它态度干扰”，A不必完全看开，可以保留自己原先的态度，只要它不会打断理论分析即可。}。
\insertparanotes
\end{itemize}
结合理论和实践，A就会得到这样的认识\paranote{同样，可以是自己琢磨出来的，也可以是听到的整套理论。}：B因为想和人说这个方面，才会去问。而之所以B不直接表达自己“想和人说这个方面”的需求，是因为某些原因：或者是出于某种含蓄的语言/社会习惯\paranote{比如说“传统文化”。}，或者是因为某种创伤\paranote{比如，之前（一些理论会认为是童年时）遇到了一个“不能表达自己需求”的环境。}，或者是因为不知道如何表达\paranote{比如说“孩子还小”，或者说“长辈以前教育条件不好”。}，或者是因为没发现自己的需求\paranote{在一些理论中，这会被解释为“一般人没有觉知自己潜意识的能力和意愿”等。}。

部分情况下，A还会做进一步的发散\paranote{比如，在引申到“长辈B试图以此来掌控A的人生”一类的结论时，会提及“这是父权的压迫”等观点。}。我们不展开讨论这些发散，而仅关注分析这一现象本身的部分。
\insertparanotes
在一些话题中，这会被当成是“发现一个人的真实动机”\paranote{心理学语境下会是“觉知潜意识”，商学语境下会是“找到真实需求”，不同的方面会带来不同的措辞。}的例子，并以此来说明“发现真实动机才能更有效地解决问题，才能从根本上解决问题”等观点。

我们容易看出，这种理论\paranote{此处不包括“只将‘反过来询问’当成是一种生活小妙招”的书籍视频播客资料，也不包括为这种现象提供其它解释的资料，而只包括“将其解释为B本身就想讲”的理论。}在将“A反过来询问B”视为一种向前竞争：“B询问A”这一行为本身由“B有讨论需求”触发，而“反过来询问B”成功地让讨论需求触发了另外的行为。
\insertparanotes
\end{explain}
\indicate{但这种归因真的是正确的吗？}
\begin{explain}
有两种不同的情况，都能产生“B经常问A，被A反过来询问后开始自己说”的现象：
\begin{itemize}
\item 一种情况是，如同上面的理论分析，B确实有“想和人说这个方面”的需求。此时，“A反过来询问B”确实是一种向前竞争。

这种情况有一个比较特殊的实际场景：A是心理咨询师\paranote{仅作为例子。实际情况中也有可能是“连线主播”“很有智慧的长者”“有一定心理知识的普通人”等，此处关注的特点是“能够自行思考现象，并且总结出理论”。}，B是来访者\paranote{仅作为例子。与A的身份相对照，B也可能是“普通人的朋友”等。}，B频繁向A询问\paranote{这个场景和前文中的描述有一定不同，比如“心理咨询师一般不会很反感”。但“B经常问A，被A反过来询问后开始自己说”的主要框架依然存在。}。A经过分析，得到了“B实际上有倾诉欲\paranote{称呼可以不同，总之指的是B“想和人说这个方面”的需求。}”的结论。A在将其告知B以后，B确认了A的猜测\paranote{很多时候，B的确认带有激烈的情绪反应，比如说“陷入回忆”、“说着说着就哭了”等。}。
\insertparanotes
\item 另一种情况是，B拥有一套“思考这个方面的问题”的行为模式\paranote{有可能是因为B本身对这方面感兴趣/比较专业，或是因为B会经常/曾经经常和别人聊这个，受到了影响。}。这一行为模式中含有多个行为，行为之间没有相互触发的顺序\paranote{比如说，可能既包含“看到一个人就上去问一问”的行为，也包括“思考其它人（可能包括B自己）这方面情况”的行为。}。此时，“A反过来询问B”，实际上是触发了另一个行为\paranote{比如说“B将自己的思考说出来”。}。此时，对于“B问A”来说，“A反问B”是一个向后竞争，而“B开始说自己的事”是一个无顺序竞争\paranote{如果我们区分得不是很细，也可以将“A反问B，B开始说自己的事”整体视为一个对“B问A”向后竞争的行为。}。
\insertparanotes
\end{itemize}
第一种情况会被用于讲解“觉知潜意识”等具体操作，不一定是针对“被人问了怎么办”事件本身的分析。而相对地，第二种情况因为更难建立理论，会较少地出现。这会导致数据来源失真，第一种情况显著多于第二种。
\end{explain}
\indicate{以上两种情况对于“B经常问A，被A反过来询问后开始自己说”不可区分。}
\begin{explain}
这个结论本身是废话。不过，一些读者（或是理论）可能会进一步认为，不仅是“两种情况不可区分”，而是“这两种情况就是同一种情况”：B所谓的“思考这方面问题的行为模式”，完全就可以看作是“想和人说这方面的需求”本身\paranote{不同的理论会对这一现象做不同的解释，比如“因为自己想了，就会想着要展示，这是自我实现的需求/肛欲期冲动/...”。我们在此不深入讨论这些理论的具体细节，仅展示这一现象。}，二者并没有任何区别。
\vspace{0.5em}
\insertparanotes
\vspace{0.5em}
\end{explain}
\indicate{将“被问了会说”的现象直接视为“有自我表达的需求”的观点，是草率且不负责的。}
\begin{explain}
主要的危险之处在于，B的“思考这方面问题的行为模式”给出的回应不稳定。这个不稳定体现在两个方面：一方面是，对于不同的B，行为模式本身会有不同；另一方面是，对于同一个B的相同行为模式，不同的刺激下，会给出的反应不同。

对于“A反问B”这一刺激，B可能的反应包括但不限于：
\begin{itemize}
\item \indicate{回答}：这就是我们上文中讨论的情况。B可能除了回答以外还会发散到别的话题，但我们都将其算作“B积极回应，现实情况按理论发展”的情况。
\item \indicate{被打断}：B没有兴趣回答A的问题，而是很快就结束了话题。这种情况可能包括“B随口应付了几句”、“B直接转移话题”、“B没有回答”等。此处还可能伴有B少量的行为，如评价“你怎么问这么多”之类的。这些行为可能很有攻击性，但只要当前话题迅速结束，我们也将其算作“被打断”一类。

这种情况已经和理论预测产生了区别。但对其应用场景来说，“B回答”和“B被打断”都能实现“A不用再受B的盘问”的效果。对于一部分A来说，因为话题更早结束，这种“不用再被B纠缠”的情况反而会更好。
\item \indicate{产生激烈反应}：严格来说，“兴致勃勃地讲自己的事”也算是激烈反应。但我们这里所列出的，是“与问题的内容无关，而对问题本身产生的反应”。

我们可以观察到几种不同的典型反应。比如说，B可能会质疑“你问这个干什么”、“小孩子别管大人的事”、“你根本就不理解我，别装着关心了”，并且将A大骂一顿\paranote{可能会发展成吵架，但后续发展对当前的讨论不重要，此处不展开讨论。}；B也可能会害怕“为什么要问我这个”、“我是不是哪里做得不好，惹到对方了”，并且情绪消沉，开始自我攻击\paranote{这种情况从外观上看和“被打断”没有明确区别，不可区分。}。

这种情况也和理论预测不同。不同于上一种情况，这种情况与理论预测有本质区别：理论本身认为“A反问B”是当前问题的解决方法，会带来一些良性的结果\paranote{主要是关系的改善和深入。“勇于沟通”一般被认为是“建立亲密关系”的正面行为。}，而实际却带来了恶性的结果，在沟通中产生了障碍。
\insertparanotes
\end{itemize}
一些理论会对后两种情况有额外的解释：B并非没有“想和人说这方面的需求”，而是“虽然有需求，但是需求被其它东西压抑住了”。比如说，“B不信任A”有可能导致B体现出“没有兴趣回复”的特点；“B感觉自己的隐私被冒犯”有可能使得B开始骂A；“B曾经被另一个自己害怕的人盘问过”有可能使得B开始害怕。
\end{explain}
\indicate{这些理论还会对后两种情况提供额外的处理方法：}
\begin{explain}
如果A想要与B建立进一步的关系，那么A就应该尝试更加深入地了解和接触B，使得B能在和A的相处中，放下恐惧和戒备，做真实的自己。这些接触方法可能包括“继续按照B的节奏，回答B的问题”、“针对性地向B展示自身的特点（比如说直接和B说‘我已经长大了’、‘你可以信任我’之类的）”等。

这些方法都有失败的可能，B可能一直都不打开心扉。但更重要的事情是，我们仍然无法区分两种情况：\indicate{B没有“直接向别人说这方面内容”，究竟是这一需求被覆盖，还是这一需求压根不存在}。即使B最终变得可以回答A的问题，我们也无法区分，B是因为“看到了自己真实的需求”，从而通过向前竞争而不再烦人地询问，还是因为“在与A的相处中，培养出了新行为”，通过无顺序覆盖而不再烦人地询问。

这两种情况都能实现“A和B之间建立了亲密关系”等目标，对它们来说没有什么区别。但是如果失败，那么“能否正确归因”就会有很关键的区别，直接决定了努力方向的正误。具体讨论参见下一节。
\end{explain}
本指南采用的态度是，只有当“B确实拥有‘想和人说这个方面’的需求”的时候，才将“B问这个方面，是因为B想跟人说这个方面”这一结论视为有效的归因。而如果B不具有这个需求，只是具有行为模式，那么即使“每次反问B的时候，B都会兴致勃勃地开始讲”，我们也不将其视作“B有需求”，认为这种情况不是“B问这个方面，是因为B想跟人说这个方面”这一结论的适用范围，用它来解释是一种外溢。

更一般地，本指南认为竞争位置是一个重要的性质。如果不考虑“某一个行为是否在向前竞争”，就会做出很多错误的归因，得到很多错误的理论。而且，这些理论的无效性不容易被发现：用这些理论指导实践时，它们经常能给出有效的方案。这是因为\hyperref[para:覆盖位置]{“在仅关注单一影响时，不同的竞争位置不可区分”}。进而，一旦方案失效，理论将既无法定位原因，也无法给出替代方案。理论的外溢反而会阻碍我们理解真实世界，阻碍我们解决现实问题。

\section{屏蔽性}
\subsection{有关现实与无关现实}
\noindent\indicate{如果一个不会自动退出的行为模式没有被打断，它就会持续下去。}
\begin{explain}
这句话听起来是废话。行为模式要不然可以自动退出，要不然不会自动退出，会无休止地反复触发一些行为。后者在不被打断的情况下，本来就可以一直持续下去。但只要换个描述角度，我们就能得到这样的结论：\indicate{行为模式如果不被打断，就可以忽略现实情况。}

这两种描述其实完全等价，但听上去的感受却有很大区别。后面的这种叙述中，我们关注的重点在于“现实情况”：一个行为模式只会响应少数特点，而会忽略绝大多数现实情况\footnote{没有数据支持，仅用于大体描述，请读者谨慎理解。关于此话题的一些展开讨论参见\hyperref[sec:行为的屏蔽性]{5.3.3小节}。}。
\end{explain}
对于一个行为链/行为模式，我们将\indicate{会触发或打断这一行为链/行为模式的现实现象}称为“和这个行为链/行为模式\indicate{有关}的现实”，或者是“这个行为链/行为模式的\source{有关现实}”。

相反，如果一个现实现象无法触发或打断这一行为链/行为模式，那么我们就将其称为“和这个行为链/行为模式\indicate{无关}的现实”，或者是“这个行为链/行为模式的\indicate{无关现实}”。
\begin{explain}
此处用到了“有关”和“无关”，但这里的定义描述的不是相关性，而是因果性。使用“有关/无关”，是因为暂时找不到更贴切的词汇，它们相比之下是最优选择。“有因果性/无因果性”过长，“有顺序/无顺序”“干涉/打扰”“有效/无效”之类的词汇指代不明。本指南中，仅当“有关现实/无关现实”整体出现时，它才表示因果性；其它地方出现的“有关/相关/无关”则仍保留原意。

虽然“忽略现实情况”的措辞看上去有贬义意味，但“有关现实/无关现实”是一种中性的判断。我们也可以见到感情色彩相反的“受到干扰/不受干扰”、“嫌麻烦/不嫌麻烦”等，对有关现实持负面态度，而对无关现实持正面态度的评价。

根据之前的定义，“无关现实”与“什么都没发生”对于这一行为模式不可区分，两个不同的无关现实对于这一行为模式也不可区分。

此处定义的“有关现实”仅涉及直接原因。比如，“现象A触发了行为B，行为B触发了行为模式C”的情况，我们仅将行为B视为行为模式C的有关现实，而将现象A视为行为模式C的无关现实。不过，如果我们将“行为B和行为模式C”整体视为一个更大的行为模式D，那么现象A就是行为模式D的有关现实。
\end{explain}
\indicate{会与行为链/行为模式A竞争，并且有可能赢得竞争}的行为/行为模式B，是行为链/行为模式A的\indicate{有关现实}。
\begin{explain}
如果B赢得了竞争，那么B就打断了A，于是B就是A的有关现实。

相对应地，如果B不会和A竞争（或者是不同时触发，或者是可以并行），或者是“虽然会参与竞争，但无法赢得竞争”，那么B就是A的无关现实。

行为/行为模式B有可能是“某种情绪/冲动”、“想完成某个目标的评估”、“某种观念/守则”等，因人而异。对同一个人的不同行为模式A，B有可能对一部分A是有关现实，另一部分则不是；对不同人的同一行为模式A，B有可能对一部分人的A是有关现实，另一部分人则不是。我们必须从实际情况出发，来分析和理解。

A和B可能没有任何其它关系，而只有“属于同一个人”和“会同时触发，并且产生竞争”这两点。而另一方面，一些和A在其它方面有关的行为/行为模式C（比如说“因为同一个现象而触发”，如“情绪感受”和“处理方法”），却反而可能是A的无关现实。

有关现实/无关现实是一种单向关系。B是A的有关现实，不说明A是B的有关现实。比如说“A和B的竞争中，总是A胜出”，那么A就是B的有关现实（因为A会打断B），但B是A的无关现实。这会方便我们研究行为/行为模式的转化链：转化链中的后一个行为模式一定是前一个行为模式的有关现实，而前一个行为模式经常是后一个行为模式的无关现实（同时也存在相反情况，比如“前者会触发后者”时，前者就是后者的有关现实）。
\end{explain}

\subsection{屏蔽与隔离}
\noindent 如果某个现象A无法打断某个行为模式B，就称行为模式B\source{屏蔽}了现象A。
\begin{explain}
行为模式B可以由现象A触发，也可以由其它因素触发。我们此处不关心触发条件。行为模式B的无关现实一定被B屏蔽。

现象A可以触发其它的行为模式C。根据定义，A被B屏蔽时，C无法打断B。很多情况下，C同时也不会触发B，此时C是B的无关现实。

我们此处讨论一些例子：
\begin{itemize}
\item 偶发的现象A：它很多时候是一个\note{盲点}。如果它被行为模式B屏蔽，那么行为模式B可能意识不到“有这么一个现象A”。于是，我们其实难以准确统计“一个行为模式到底屏蔽了多少现象”，越是偶发，就越是难以意识到它的存在。

\item 稳定存在的现象A：可以指“每时每刻都出现”，也可以指“在特定环境下以高频率出现”，我们这里主要关注“会重复提供相同的刺激”这一特点。行为模式B有可能一直意识不到现象A，也有可能“虽然能意识到（此处多指某个负面影响），但无法控制自己”。

\item 预期出现的结果A：在做计划时，对某一行为/行为模式B所带来后果的估计。此时的“现象”指的不是“结果”（因为结果还未发生），而是“预期”（作为一个认知而存在的客观现象）。如果没有对预期的认识，那么即时确实有负面现象，也不算做“被阻断”。预期不一定准确，但我们此处不关心其正确性，只关心是否可以打断。

\item 需求/困难/目标A：这个语境下，A是一大类有所不同的现象，而“打断”指的是“因为失败而换方法”。行为模式B在拥有“失败”这一终止条件（并且同时有对应的判断方法）的情况下，能屏蔽一类现象A，意味着能够解决A中的所有不同情况。
\end{itemize}
“屏蔽”一词主要指“现象A无法对行为模式B的影响向前竞争”。从语感上，这一概念屏蔽了“现象A”对“行为模式B的影响”的影响。我们此处要求的不是“现象A无法影响到行为模式B的影响”，它仍然可以向后竞争；而是“现象A通过影响行为模式B，从而影响行为模式B的影响”这一条影响途径被阻断。

当“行为模式B的影响”不是那么容易改变时（比如说“态度”、“心情”等），“无法向前竞争”在很多情况下就无法带来任何改变。此时，现象A是“行为模式B的影响”的无关现实，行为模式B确实完全隔离了“行为模式B的影响”与“现象A”。
\end{explain}
如果当行为（行为模式）B不触发时，某个现象A会触发某个行为（行为模式）C，但B触发时，A就不会触发C，那么就称行为（行为模式）B\source{隔离}了现象A与行为（行为模式）C。
\begin{explain}
“屏蔽”与“隔离”描述的现象虽然相近，但也有明显的不同。比如，B屏蔽现象A时，不一定隔离了现象A和“B的影响”，不一定就会使某个行为C不触发。而“隔离”也没办法很好描述“只有现象A不会触发C，而B触发时C也触发”的情况。为了避免术语过多，此处我们也不再引入“导通”之类的概念。

如果我们不是很关心时序和具体的行为触发链，我们可以将其视为“现象A作为刺激，行为B和行为C产生竞争，B覆盖了C”。当我们落实到具体的行为触发链时，情况可能更复杂一些：A有可能不直接触发C，而是先触发了数个行为后，再触发C。B可以在此过程中与其中的任何一个行为竞争并取胜，打断这一行为链，并且使得C不被触发。

我们此处讨论一些例子：
\begin{itemize}
\item 转移注意力：现象A先出现，并且触发一个行为链，这一行为链会触发C；在此过程中，现象D出现，触发行为/行为模式B，打断了这一行为链，使得行为C不再触发。比如，在自己没事做（现象A）的时候，感到悲伤、孤独、被抛弃（行为链）。如果发展下去，就会深深地emo\paranote{此处我们不关注这是病理性抑郁，还是网抑云听多了，或是其它情况。具体情况的不同不影响此处讨论。}（行为C）。这时候如果有人找自己聊天（现象D），聊起来感觉好受很多（行为模式B）。
\item 先入为主：行为模式B先出现\paranote{可能由某个现象D触发，或是独立存在。}，现象A后出现，并持续存在。现象A原本会触发行为模式C，但是行为模式B反复参与竞争，导致无法实现C。比如，到了学习/工作的时间（现象A），但是有其它的干扰因素\paranote{包括但不限于“想玩”、“讨厌学习”、“（其它原因的）心情不好”、“周围环境太吵”、“老是有人找自己聊天”等。具体情况的不同不影响此处讨论。}（行为模式B）。在无干扰状态下可以顺利进入状态，完成任务（行为模式C），但在干扰下，虽然人坐在桌子前，但是总是进不去状态，也没有任何进展。
\item 改主意：现象A触发了行为模式C，行为模式C（直接或间接地）触发了行为/行为模式B，之后B覆盖了C。比如，自己第一次见到一个看上去很吓人的东西\paranote{可能是某种操作流程，可能是某个人/生物，可能是某种（很难的）知识。}，觉得自己肯定接受/处理不了，但冷静下来想了想，又发现自己好像能接受/处理。我们之前介绍过的“通过向后竞争，得到用于向前竞争的行为”，就是这种情况。
\end{itemize}
\insertparanotes
当我们关注行为（行为模式）C时，行为（行为模式）B在对C向前竞争，并且成功地向前覆盖了C。从上面的例子可以看出，向前竞争的B有正面的也有负面的，其判断主要取决于C本身是正面的还是负面的。此时，如果不去和B竞争，而只是强调“看到A就应该做C”，就不会起到效果。
\end{explain}

\labelsubsection{行为的屏蔽性}
\hfill\begin{minipage}{0.6\textwidth}
\fontsize{8pt}{12pt}\selectfont\fontsize{8pt}{12pt}

\raggedright 听不见胸腔中绝望的声音，却呕出了鲜红的悲戚。\footnote{\bilibili{BV1j1421k7Zg}。\\}

\raggedleft JUSF周存《公主》
\end{minipage}

\noindent\indicate{原子行为会屏蔽所有现象。}
\begin{explain}
这句话看起来有些奇怪，因为定义\hyperref[def:屏蔽]{屏蔽}的时候没有定义过“行为的屏蔽”。硬套定义也不可行，因为行为不存在“打断”这一概念，只有“触发”和“没有触发”。另一方面，如上一段所说，此处的“屏蔽所有现象”不包括“突发的生理性损伤”等，而只关注意识现象。

这个结论有两个作用：
\begin{itemize}
\item 一个是作为补充定义。如果一个行为模式可以拆成多个部分，那么“这一不屏蔽的现象”就能分成“每个部分不屏蔽的现象”和“这些部分触发时，参与竞争（并可能占优）的现象”两种来源。将“原子行为不屏蔽的现象”设为“无”，就能让这个拆分过程走到最后，将“不屏蔽的现象”完全归于竞争过程之中。

此处需要注意，我们不将“原子行为受到的刺激”视作“原子行为不屏蔽的现象”，因为此时原子过程还未触发。在刺激引起的竞争过程之后，这一行为才会执行，此时已经不会再受刺激了。

\item 另一个是区分原子行为与复合行为。复合行为之中包含行为触发，也就可能有竞争过程，以及对应的不屏蔽现象。这些竞争过程虽然不一定存在，现象也不一定真的能赢得竞争，而且它们也不一定可见，但这是一个不可或缺的理论工具。

比如，当刺激A客观存在，且不可消除时，如果我们想要改变一个原子行为的结果C，那么只能向后竞争\paranote{此处可能出现“随着逐渐熟练，原子行为和后续行为逐渐形成了新的复合行为，进而演化出新的原子行为，覆盖原有原子行为”的现象，不过不是此处讨论重点。}。而对于一个复合行为，比如“A触发了B，B又触发了C”，更多的触发带来了更多的竞争位置，我们就也可以对C向前竞争。向前竞争的手段有可能用于消除B的影响，这样的手段对原子行为没有针对性的效果\paranote{在一些情况下，仍然可以和原子行为向后竞争，进而覆盖。但这不具有针对性。}。

再比如，我们可以得到如下结论：\indicate{一个原子行为形成后，就会屏蔽它的形成环境。}\paranote{一部分读者可能会将其视作“在作品完成之际,作者就已经死亡”一类的观点，而另一部分读者可能会将其视作“公理系统不关心自身的模型如何构建”之类的观点。它们虽然有相似之处，但语境和侧重点有较大差异，不建议读者直接套用。}此处的“形成环境”可能是“它由某个复合行为简化而来”\paranote{由于人（以及其它动物）真的可以习得条件反射，我们必须允许“形成一个原子行为”的机制，而不能将其简单视为“只是某个熟练的复合行为”。}，也可能是“从外部环境中习得”，没有什么限制。如果从复合行为中简化而来，那么复合行为有可能会并行触发，但打断复合行为无法起到覆盖原子行为的作用。

这一结论远比“原子行为会屏蔽所有现象”弱，但它在一些情况下会更好用，比如：如果我们想要覆盖一个原子行为，那么对“这一原子行为如何形成”的任何归因都无法起到针对性作用，任何的理论都只能参与向后竞争。能否成功不由理论的正确性决定，而由理论的说服力决定\paranote{此处的“说服力”并不取决于“理论多有道理”，甚至不总是取决于“理论听起来多有道理”，很多时候主要由“理论的来源”等和理论内容没有直接关系的现象而决定。只要能竞争成功，让人认同，那么理论就“有说服力”；反之则“没有说服力”。}。
\end{itemize}
\insertparanotes
\end{explain}
\noindent\indicate{一个多重触发的行为，会屏蔽其所有触发途径屏蔽的现象。}
\begin{explain}
更严格的叙述如下：一个多重触发行为会屏蔽的现象，是其所有触发途径所屏蔽现象的并集；而一个多重触发行为不屏蔽的现象，是其所有触发途径不屏蔽的现象的交集。

对于\trained，本结论和上一段的论述有更易于理解的描述方式：串联的时候，一处断路就是断路；并联的时候，一处通路就是通路。读者也可以据此研究触发关系更复杂的行为的屏蔽情况。

如果我们想要用一个新行为来覆盖多重触发的行为，那么这个新行为很难对每个触发途径都向前竞争。大体上来说，触发途径越多，则越难完全地向前竞争\paranote{单纯从“数量”上来衡量，则会有“无论有多少个触发途径，都能通过消除刺激的方法向前竞争”和“触发途径只有一个原子行为，但刺激不可消除，只能向后竞争”的反例。此处的严格表述为，我们对多重触发的行为以“包含”建立一个偏序关系，那么如果一个行为的触发途径完全包含于另一个，那么这个行为向后竞争的触发途径也完全包含于另一个（于是数量自然不多于另一个）。}。在部分情况下，这一新行为可以通过向后竞争来完全覆盖多重触发的行为；但另一些情况下则不能：我们经常能见到“虽然诚恳且深入地意识到了自己的错误，但还是无法控制自己”等情况，这一情况经常是因为“认识到错误仅能打断其中一个触发途径，但还有别的触发途径”。

另一种覆盖多重触发的行为的思路是“为每个触发途径单独选择可以向前竞争的行为”，这需要对这一多重触发的行为有全面而准确的理解，漏掉任何一个触发途径都有可能导致无法覆盖。
\insertparanotes
\end{explain}

\labelsection{能力的形成}
\subsection{行为的形成}
% 对于一个行为和一个环境，我们将\indicate{在该环境中，稳定地触发一个与该行为等价的广义行为}称为这一行为（在该环境中）的\indicate{形成}，广义行为称作这一形成过程的\indicate{行为雏形}或是这一行为的\indicate{雏形}，该行为称作这一形成过程或是该广义过程的\indicate{目标行为}。
% \begin{explain}
% 定义中提到的“行为”仅根据“刺激”与“行动”定义，并不要求是“某个具体的人拥有的原子行为”。定义中提到环境，主要是为了控制刺激与竞争相同，具体分析见后文。

% 我们在3.4.2小节中定义过\note{关联}的建立。“行为的形成”与“关联的建立”的主要区别在于，关联的建立只计原子行为，而行为的形成也计复合行为。
% \end{explain}
如果\indicate{几个（原子）行为通过数次合并，形成了一个新（原子）行为}，那么我们就将这几个行为称作这个新行为的（一个）\indicate{雏形}，将雏形和这些合并合称为这个新行为的（一个）\indicate{形成过程}，将合并时对应的竞争称为这一形成过程\indicate{包含}的竞争。

在一个环境中，如果一个广义行为中的行为是由某一些行为合并产生，那么就将这些行为称作这一广义行为的\indicate{雏形}，将\indicate{雏形、这些合并、上一个行为和下一个行为的连接}三者合称为这个广义行为的（一个）\indicate{形成过程}/\indicate{组装过程}，将合并时与连接时对应的竞争称为这一形成过程\indicate{包含}的竞争。

我们将\indicate{新（原子或广义）行为}称为这一形成过程的\indicate{主要行为}。对于形成过程包含的某个竞争，我们将“形成过程中，赢得竞争的行为”称为该竞争的\indicate{目标行为}或是该形成过程的一个\indicate{目标行为}；将“参与这一竞争，但不是目标行为的行为”称为称为该竞争的\indicate{非目标行为}或是该形成过程的一个\indicate{非目标行为}。
\begin{explain}
依照这一定义，同一行为可能有多个不同的雏形。在原子行为的“雏形”定义中，我们允许“原子行为进行零次合并”的情况，即“将原子行为本身（组成的一元集合）也视为自己的雏形”，并将其称为平凡雏形/平凡形成过程，其它雏形则称为非平凡雏形/非平凡形成过程。这样处理会方便行为链的“雏形”/“形成过程”定义。

虽然我们使用了“形成过程包含的竞争”的称呼，但形成过程中并不直接包含竞争，而只包含“对应的行为赢得了竞争”的现象。具体的竞争会根据环境的不同而不同，但只要占优的行为相同，我们就将其视为同一个形成过程。%参与竞争但被覆盖的其它行为都是该行为的无关现实。
因此，我们可以认为，形成过程包含所有的目标行为。但相反，它不包含所有的非目标行为\footnote{对于\rigorous，此处需要排除“某个竞争的非目标行为是另一个竞争的目标行为”的特殊情况。}。

“找到一个原子行为的非平凡雏形以及对应的形成过程”就是\note{理解}了这一原子行为（即“给该原子行为溯因”），相关的讨论参见\hyperref[sec:对竞争过程的解读]{4.4.3小节}。

广义行为的“雏形”和“形成过程”定义中，不要求这些子行为是原子行为。它们也可以是广义行为，此时我们不递归定义（即，不地将子行为的“雏形”和“形成过程”也视作原行为的“雏形”和“形成过程”的一部分），而只是将子行为视为雏形中的一个元素。这种处理可以方便我们后续的理论分析。

广义行为的“形成过程”定义中，“上一个行为和下一个行为的连接”指“上一行为触发下一行为”的过程。广义行为的“雏形”/“形成过程”定义后半段的严格表述是：广义行为的雏形，是“其中每个行为的雏形”的并集；广义行为的形成过程，是“其中每个行为的形成过程”，与“所有连接”的并集。同样地，同一个广义行为也可以有多个不同深度的雏形/形成过程。

我们不要求“（原子/广义）行为的雏形”也是一个广义行为，更不要求“它是一个和主要行为等价的广义行为”。雏形可以是一堆杂乱无章的行为，它们之间可能没有触发顺序，或触发顺序与主要行为相反，或是部分相同部分相反部分无关。这能方便我们将更多的现象纳入分析，详细的讨论参见\hyperref[sec:行为的转化与合并]{3.4.3小节}。如果雏形也构成一个与主要行为等价的广义行为，就也可以沿用“行为链\indicate{简化}成行为”的称呼，将这种现象也称为“雏形简化成了原子/广义行为”。不过我们此处不正式引入这一术语。

与原子行为不同的是，广义行为的雏形/形成过程额外要求“在同一环境中”的条件。这一限制要求“这些合并必须在该环境下产生”，以避免某些过远的溯因。这使得“什么可以算广义行为的雏形/形成过程”与环境的选择相关。

不是所有的行为都有非平凡雏形，我们也将“直接从外界习得行为”视为一种独立的方式，这种行为仅有平凡雏形。在一个环境中，“无法进一步溯因”的原子行为，和“直接从外界习得”的原子行为不可区分。我们也可以因此将“（在某个环境中）从外部习得”更清晰地定义为“（在某个环境中）无法进一步溯因”，这使得“从外部习得”明确地和环境（视角）相关。
\end{explain}
\indicate{使用某个形成过程形成（原子/广义）行为}的充要条件是\indicate{具备该形成过程对应的雏形，并且在形成过程包含的每个竞争（包括合并时和连接时）中，对应的行为赢得竞争}。
\begin{explain}
此结论仅针对某一具体的形成过程，而不关注其它形成过程，以及“行为是否可以由其它形成过程形成”。相似的讨论在\hyperref[sec:有效沟通的必要篇幅]{4.2节}中也有涉及，此处不再重复。

此结论本身是显然的，结论中的两个要求（具备雏形和赢得竞争）本身就是形成过程的定义。此处将这一结论单独提出，是为了方便后续的讨论。\trained 也可以将其视为“判断两个不同行为的形成过程等价”的条件。

结论中并未直接涉及到“环境”，相同的形成过程也确实可以在不同的环境中发生。环境在形成过程中的作用是“确定参与竞争的行为”。行为的不同有可能来自于刺激的不同，有可能是来自于个人已有行为的不同。
\end{explain}

\subsection{主动形成与被动形成}
对于一个行为和一个形成过程，一个行为模式/人A和一个人B来说，如果A能够判断B具有（形成过程对应的）雏形中的所有行为，并且能使对应的行为赢得对应的竞争，那么就将“B按照这一形成过程形成这一行为”称为“A使B按该形成过程\indicate{主动形成}了这一行为”，或是“A按该形成过程\indicate{培养}了B的这一行为”。
\begin{explain}
行为模式A不一定属于人B，它可以是B拥有的行为模式，也可以是另一个人拥有的行为模式。B在此处的作用是“提供雏形中所需的行为（以及这些行为对应的覆盖关系）”，这些行为不一定需要“全部属于某个行为模式”，也不需要“能组成一个广义行为”。

我们沿用上一小节的定义，将“要培养的行为”和“雏形中的行为”分别称作“主要行为”和“目标行为”。根据B的个人情况不同以及外在环境不同，“能使目标行为赢得对应的竞争”需要做的事、难度、以及是否可行，都会有所区别：
\begin{itemize}
\item 理想状态下，B具备目标行为，且目标行为本来就可以赢得竞争。此时不需要做任何额外的工作。只要环境具备\paranote{比如说一份详细的教程，或者一个合适的比方。}，B就能自行组装出主要行为。
\item 如果B具备目标行为，但目标行为无法赢得竞争\paranote{包括但不限于“会参与竞争，但是会被覆盖”、“会被其它行为隔离，无法参与竞争”等。}，那么就需要通过其它手段来处理这一竞争中的其它行为。通过向前竞争的手段\paranote{直接使其它行为不触发（从而不参与竞争），或是无法占优。}或是通过向后竞争的手段\paranote{比如形成一个“从赢得竞争的行为到目标行为”的行为链。这个行为链本身可以视为一个“和该目标行为等价”的复合行为。}都可以解决问题，具体采用哪种方式，依实际情况而定。

需要注意的是，无论是向前竞争还是向后竞争，都会涉及到“培养新行为”。这并入下一种情况讨论。
\item 如果B不具备目标行为，那么就需要使B形成这一目标行为，即“对这一目标行为递归地使用上述理论”。培养这一目标行为可以不限制具体的行动过程。当然，也可以考虑“更换主要行为的形成过程”，使得不需要这一目标行为即可形成主要行为。
\end{itemize}
\insertparanotes
我们在定义“行为的主动形成”时，不递归地要求“如果过程中涉及了额外的行为形成，那么这一行为也需要是主动形成的”，只关心“雏形中的行为”和“这些行为相关的竞争”两方面。
\end{explain}
我们将\indicate{B形成了一个行为，但这不是A使B（按任意形成过程）主动形成的行为}的现象称为“B（对于A）\indicate{被动形成}/\indicate{巧合形成}了这一行为”。
\begin{explain}
此处的“巧合”和\hyperref[sec:巧合有效的沟通]{4.3节}中的“巧合”含义一致。

如果A知晓多个不同的形成过程，并且B按照其中的一个形成过程形成了行为，那么我们仍然将其视为主动形成。此处定义的被动形成，主要针对的情况是“A不知道B这一行为的形成过程”。

讨论“行为的被动形成”，和讨论“所有的行为形成现象”没有区别，也和讨论“普遍的竞争-关联机制”没有区别。竞争-关联机制的复杂性主要来源于三方面\paranote{其中前两方面是\hyperref[para:竞争的随机性]{3.4.2小节}中已有的讨论内容，主要关注单一竞争的复杂性。}：
\begin{itemize}
\item \indicate{参与同一个竞争的行为不同}：即，对于相同的刺激，和参与竞争的某一相同行为来说，“参与竞争的其它行为”与“该行为与其它行为的覆盖关系”是未知信息。比如，“某个行为在某一竞争中获胜”（如“成功经验”或是“对情况的预测”），并不能用于推导“这一行为会在刺激相同的其它竞争中获胜”（即“照着这么做就好了”）等结论。
\item \indicate{竞争会受到环境的明显影响}：即，竞争经常不是单独出现，而是处于一个行为链中，每个竞争的参与者都由上一个（或者多个）竞争的胜者决定。这造成了“即使行为相同，它们竞争的顺序不同，就会有不同的结果”、“不处于某行为模式时，行为A可以赢得竞争；处于某行为模式时，则是行为B赢得竞争”等现象。
\item \indicate{刺激与其触发的行为不可控}：人会同时接触到很多现象，每个现象都有可能触发行为；每个现象都具有多方面特征，每个特征都有可能触发行为。由此可能产生“同时有多组竞争在并行”、“实际出现的竞争不是由正确的刺激（指‘某个形成过程中的刺激’）触发”等现象，同一个环境中演化出的行为可能有显著不同。
\end{itemize}
对于某一个形成过程来说，有很多种被动形成的行为需要注意，如：
\begin{itemize}
\item \indicate{占优的非目标行为}：作为形成过程的有关现实，需要通过向前/向后竞争处理。相反，无法赢得（任何一个）目标竞争的行为，就是形成过程的无关现实，不需要关心。
\item \indicate{未知且占优的非目标行为}：从属于上一种情况。此处的“未知”主要指的是“无法确认触发条件”，或是在复合行为中“无法确认某一步骤的效果，即下一步骤的触发条件”，从而无法向前竞争，只能向后竞争。手段受限经常会导致无法成功覆盖。
\item \indicate{新出现的被动形成的行为}：由“形成过程本身的竞争”、“覆盖非目标行为时的竞争”等现象而产生的，与目标行为不同的行为。它可能在“触发条件”“竞争能力”“实际行动”等方面与目标行为有差别。这会导致“培养出了一个新行为A\paranote{A可能是某个目标行为，也可能不是。}，但这个A会参与另一个竞争，并且覆盖那个竞争的目标行为B，使得B更难以形成，或无法形成；形成B后又会反过来覆盖A”等现象。
\end{itemize}
\insertparanotes
\end{explain}

\subsection{行为模式的形成}
我们将\indicate{一个人接触某些现象时，形成的所有由这些现象触发的（广义）行为}称为“这个人因接触这些现象而\indicate{形成}的行为模式”，或“（这个人的）（对这些现象的）本次接触\indicate{形成}的行为模式”。
\begin{explain}
定义中包含“由这些现象触发”的要求，是因为在接触这些现象的同时，还可能会有其它的并行行为。这些行为自然不算。当然，如果这些并行行为和现象带来的刺激产生了关联，那么就会算入“因接触这些现象而形成的行为模式”之中。

“接触”某些现象，指的是某次具体的接触，而不只是“和这些现象的所有接触”。当然，“具体的接触”中也包括“今天接触了一次，明天又接触了一次”、“这几个月以来，这些现象一直在身边不断出现”之类连续发生的多次接触，所以我们不会漏掉什么情况。
\begin{itemize}
\item “接触某些现象”不一定是首次接触（或是“从首次接触开始的接触”），“某次接触时形成的行为”也不一定就是“这些现象会触发的行为模式”，后者可能包含更多行为。但当我们如果条件取为“在当前的拥有的行为和其组织形式下，接触这一现象”，那么“本次接触形成的行为模式”就可以视作一个行为模式\paranote{但读者也可以将其视为对“行为模式”这一术语的滥用。本指南倾向于尽量不要引入过多新概念。在实际行文中，读者应能自行分辨两种不同含义的“行为模式”。}。
\item 如果拥有的行为和组织形式发生了改变（比如说“十年前接触了一次，十年后改变了很多，又接触了一次”；或是“两次接触之间发生了重要的事情，心态和行为产生了明显变化”），导致“上一次接触时的行为”因为其它原因产生了额外的演化，那么这就不能视作一个行为模式，我们也不将这种“有间断的多次接触”视为一次连续的接触。
\end{itemize}
\insertparanotes
此处的“某些现象”可以指代的内容有：
\begin{itemize}
\item \indicate{某个（稳定存在的）环境}：此处我们主要关注的是“（在某个行为模式存在的时间范围内）稳定存在，不会自行消失”的所有现象，不局限于“某一个现实存在的环境”。

我们可以关注“与环境的第一次接触”，以研究由此形成的初见印象；也可以关注“从第一次接触到现在的所有接触”，以研究稳定下来的行为模式；或是关注“某个时间点（一般是行为剧烈变化的时间）附近的接触”，以研究该时间点附近的行为演化。
\item \indicate{某个一次性事件}：这里指的不是“除了稳定存在的环境以外的所有现象”，因为还有另一些“会多次重复出现，但会在一段时间后消失”的现象。我们主要关注那些“会给人带来深刻影响”的一次性事件，它们的特征比较突出，可以用来对比。

常见的此类现象包括但不限于“某个突发的重大消息”、“记了一辈子的某一句话”、“某段不敢回忆的被封闭的经历”等等。这可以说明，现象对人的影响不由“接触的时长”而决定，而由“深入内心的程度”（这只是“形成/覆盖了多少行为”的一种更感性的描述）而决定。
\item \indicate{某个评价或要求}：从现象的量来说，这更接近于“一次性事件”，因为它只有一句话（一个观点）；从接触次数来说，这更接近于“稳定存在的现象”，但也不排除“一句话记了一辈子”的情况。

A评价B时（A和B可以是同一个人或不同的人），B有可能会因为这些评价而形成一些行为。这些行为有可能对A来说是主动形成或被动形成，有可能在或不在A的预期之内，有可能是A所期望或不期望的，有可能对A产生某些影响。
\item \indicate{形成行为的尝试}：无论是“想要形成一个行为”的念头本身（或是“有别人说要这么做”），还是为此而做的所有行为（如收集和接触资源、尝试、调研、训练等）和接触到的所有外部刺激（如资源、教育、评价、与人的冲突等），都可以视作是此类的现象。

我们暂时将想要形成的行为称为目标行为。接触这些现象，不代表可以形成目标行为。实际的竞争、关联、行为形成可能不符合预期。如果对应的刺激与另外的行动关联为了非目标行为，那么它的存在通常会使得目标行为更加难以形成（因为它会在刺激产生时参与竞争，使得目标行为更加难以胜出）。
\end{itemize}
\end{explain}

\labelsubsection{多行为的形成}
\noindent 我们将\indicate{多个（原子/广义）行为组成的集合}称为一个\indicate{行为组}，将\indicate{这个行为组中，所有行为的触发条件组成的集合}称为它的\indicate{刺激组}。
\begin{explain}
在这里才定义行为组有点过晚了，主要是因为之前想着能否避开这个定义。使用这个措辞的话，之前的很多定义（比如“行为的雏形”）都能得到简化。

之前也提到过很多具有其它性质的行为组，比如“在同一个人身上，同时触发，但无法同时执行”的多个行为，就是一个“会引起竞争”的行为组。

再比如，对于行为链或是行为模式，如果我们不关心其中的行为触发与覆盖关系，而只关心这些行为本身，那么它们也可以视作一个行为组。这可以方便我们稍后讨论“形成行为模式”时的具体处理。
\end{explain}
我们将一个\indicate{至少包括两个行为，且这些行为的触发条件和执行的行动都不相同}的行为组，那么我们就将它们称为一个\indicate{需求信息}的行为组。
\begin{explain}
这一定义具体指的是“对于行为组中的任意两个行为，它们的触发条件不相同，且它们执行的行动不相同”。

% 其实这个定义可以收紧一些：如果行为组中，有多个触发条件相同的行为，那么我们将它们整体视为同一个行为，其行动视为“所有这些行为的行动组成的集合”（而如果某个刺激只对应一种行动，那么就将其行动改为“只包含这一行动的集合”）。然后，再将“行动相同的行为”的刺激合并（合并方式与合并行动时一致）。此时，如果还剩至少两个行为，且这些行为两两之间满足“刺激和行动均不相同”，那么就将其也称为“需求信息的行为组”。这可以用于“刺激A对应行动1，刺激2对应行动1和行动2”之类的情况。但本指南在后续分析中，仍采取现在这种较为方便的定义。
其实这个定义可以收紧一些：如果行为组不能表示为“一个非空刺激集合和一个非空行动集合的笛卡尔积”（即，不是“对于某些情况，具体怎么做都行”的行为组），那么就将其也称为“需求信息的行为组”。这种定义包含“刺激A对应行动1，刺激2对应行动1和行动2”之类更复杂的行为组。但本指南在后续分析中，仍采取现在这种较为方便的定义。读者可以自行将这些分析扩展到更一般的情况。

此处的“需求信息”主要指的是：有“区分不同情况”的必要性。如果行为的触发条件相同，那么就没有“不同的情况”需要区分；如果行为的行动相同，那么“区分情况”是这一行动的无关现实。只有在触发和行动都不同的时候，才需要“根据不同情况做出不同的行动”。

行为模式和行为链经常可以视为“需求信息的行为组”，因为其中的每个行为的触发条件是“上一个行为”，这些行为两两不同的时候，它们就满足“触发条件和执行的行动都不相同”的情况。但行为模式也有“无论发生什么，都会触发某几个固定的行为，然后这些行为之间相互触发”等情况，这种行为模式就不需求信息。
\end{explain}
对于一个需求信息的行为组和一个视角，如果\indicate{这一视角无法区分这一行为组的刺激组}，那么就称这一视角对这一行为组有\source{信息缺失}、\indicate{信息需求}或\indicate{调研需求}。
\begin{explain}
为定义简便，此处将“刺激”直接视为现象。对于\rigorous，可以先通过“一个现象能触发行为组中的一部分行为，而无法触发其它行为”来定义“这一现象对行为组有区分性”，再根据区分性定义“需求信息的行为组”，此时就有了“区分此行为组的一些现象”，之后再判断“视角能否区分这些现象”，以此完成上述定义。

“行为组的触发条件”（即“可以触发行为组中的哪些行为”）形成了一个视角B。如果一个视角A对这些行为组有信息缺失，那么视角B就可以为视角A提供（区分这些行为的行动的）\note{必要信息}。这也是“信息缺失”这一措辞的由来。当然，视角A不一定要从视角B这里获得必要信息，视角B甚至不一定在任何人身上存在。
\end{explain}
\indicate{如果一个人（行为模式）对一个行为组有信息缺失，那么这个人（行为模式）无法同时形成这个行为组中的多个行为。}
\begin{explain}
此处是“无法同时形成”，指的不是“可以先形成其中一个，再形成另外一个”，而是“如果形成了新的，那么旧的会被覆盖”。而“形成了其中一个”也指的是“在对应的刺激出现时，能做出对应的行动”，仅保证行动相同，而触发条件一定不同。

如果一个人身上有多个信息缺失的行为模式，且这些行为模式之间会相互屏蔽（可以并行，或者是“虽然同时只能有一个触发，但是触发时不会打断其它行为模式”），那么不同的行为模式可以形成不同的行为。但是这些行为仅有行动部分与行为组中的某个行为一致，触发条件则一定会不同。此时我们应该将每个行为模式分开考虑，仍然符合这一结论；而不应该将这种现象视为“一个人同时形成了这个行为组中的多个行为”。

这一结论本身是显然的：一个行为模式有信息缺失，就说明“实际的触发条件”不存在这一行为模式中（有可能是不具备接触条件，也有可能是被其它行为\note{隔离}了），从而它不可能与对应的行动合并成准确的行为。实际形成的行为的触发条件，一定是某个“无法区分这些行为的刺激组”的条件，从而会响应所有的刺激，于是产生覆盖。

我们设上述形成的行为的刺激为“刺激O”，而行动是“行为组里面的行为A”的行动。我们将其简记为“行动A”，而将行为A的触发条件记为“刺激A”。其它行为的名称依此类推。这一行为在出现刺激A时可以正确响应，但出现刺激B时则不行。此时如果想要将“刺激B时的行为”替换为行动B，则最终的竞争结果是“刺激O与行动B合并，形成了新行为”。这个时候，刺激A的响应又不对了。

这种情况有两种解决方法：\vspace{-0.25em}
\begin{itemize}
\item 一种是向前竞争，让刺激O不出现。这在“刺激O（以及对应的行为）隔离了有效信息”的情况下是可行的方案，但如果行为模式中的其它可选刺激也都无法区分，则会无效。\vspace{-0.5em}

\item 另一种是向后竞争，让刺激O与某个“可以得到更多信息的行为/行为模式”合并。与上一种情况类似，如果有可选的行为，那么就可以这么做；如果没有则不行。\vspace{-0.25em}
\end{itemize}
如果这样可以得到有效信息，那么就可以分别形成这一行为组中的行为。而如果没有可以提供信息的行为，就需要先形成这一行为。但这样就会遇到另一个问题：
\end{explain}
\indicate{形成行为时，如果需要递归，则必然存在信息缺失。}
\begin{explain}
“递归”是指的“为了形成主要行为，需要先形成某一目标行为，以满足雏形的需求”的过程。信息缺失指的不是“主要行为”和“目标行为”之间需求信息，毕竟主要行为就来自于目标行为的合并，它们的刺激或者条件经常相同；而是指的“形成主要行为”和“形成目标行为”之间需求信息：

二者的刺激（动机）不同，前者是任意的（但绝对不是“形成主要行为”），而后者是“形成主要行为”；二者的行动不同，一个是“主要行为”而另一个是“目标行为”。二者都是一次性的行为，需要当场形成\footnote{此处的叙述不严谨，没有经过完整的分类讨论，仅供读者大致理解。}，所以很容易搞乱。

避免的方法也有两种：一种是“一次只形成一个行为”，比如“先打好基础再考虑进阶的难题”之类的，在形成目标行为之前，不要考虑形成主要行为的事情。这在理想的引导下可以实现。另一种是“获取对应的信息”，也就是“思路清晰地意识到每一步的作用”，即“主动形成一个行为”。这需要对行为的形成有较为深入的认识。认识得越准确，则越能驾驭复杂（指“多层递归”）的行为形成流程。
\end{explain}

\section{总结与讨论}
\subsection{本章总结}
本章内容围绕\indicate{行为形成}的过程和原理而展开，讨论“行为的形成过程”，以及这一现象所需要注意的方面。

我们将行为的形成视为一个递归的过程，即“一些行为（作为\indicate{雏形}），通过某些竞争（作为\indicate{形成过程}），形成新行为”。严格来说，递归性不是“形成行为”的性质，而是“由行为形成行为”的性质。当行为\indicate{被动形成}时，只需要产生“某一个竞争”，此时不需要考虑递归性。只有当我们想要\indicate{主动形成}一个行为时，才需要去顺着因果链条去递归地研究，反向设计一个可行的形成过程。

“雏形”不足以完全刻画行为的形成，即，不是“只要雏形中含有对应的刺激和行动，就能形成行为”。具体能形成什么行为，由竞争过程决定。但是，相比起行为，竞争过程在很多时候是不可知的。为了刻画这种现象，我们将其抽象，引入了基于\indicate{视角}的\indicate{不可区分性}。

我们主要关心不可区分性的两个方面：一方面是“不可区分会带来信息缺失”，这可以用来刻画“被覆盖的行为和不存在的行为不可区分”（也即“竞争过程不可知”）等现象；另一方面是“不同现象的区别在一些时候无关紧要”，这可以用来刻画“\indicate{原子行为}和\indicate{复合行为}都能参与行为形成，可以相互替换”等现象。这两个方面虽然在作用上有明显不同，但它们是相同现象的不同表现，所以在行文中没有特别区分。

处理竞争过程，也就是要“触发目标行为，打断/覆盖/控制非目标行为”。这要求我们去具体地研究“每一个行为会在什么时候被触发，又会在什么时候被覆盖”。因此，我们引入了“\indicate{有关现实}/\indicate{无关现实}”的术语，来描述那些“能产生影响/不能产生影响”的事情。此处的“有关/无关”可以按照“这一信息能否应用到行为形成”上来理解。作为覆盖机制的细化，我们又引入了\indicate{屏蔽}和\indicate{隔离}两个概念，用于判断“使用哪些现象是有效的，哪些现象是无效的”。

准备好了理论工具后，我们就可以来分析和处理实际的竞争过程了。作为一个简单易懂但十分重要的例子，\indicate{原子行为}和\indicate{复合行为}的屏蔽性有显著差别：原子行为无法中途打断，而复合行为则可以中途打断。我们将这种区别抽象出来，总结成了\indicate{竞争位置}的概念：复合行为可以\indicate{向前竞争}，而原子行为只能\indicate{向后竞争}。更多的竞争位置为我们提供了更多的打断/覆盖机会，更容易让竞争向着所需的结果发展。

与此同时，“竞争位置”还是一种良好的分析工具。它可以帮助我们定位不同的竞争，以便更准确地解读和改变意识现象。“参与了不同的竞争，从而起到相同效果”的现象广泛存在，如果我们将其笼统地视为同一种方法，并且生搬硬套到其它情况上，就有可能产生预期之外的严重后果。这种\indicate{局部有效理论}忽略了实际存在的复杂性，试图在有\indicate{信息缺失}的情况下达成自己的目标。这不可能稳定成功。

\labelsubsection{本指南的分析思路}
可能会有一些读者觉得本指南的讨论戛然而止，只是罗列了可能存在的问题，而没有针对这些问题给出具体的解决方法。在看完第五章后，仍然不知道应该如何形成新行为，如何成长。这种感觉是正常且合理的。这主要有以下四个原因：\vspace{-0.25em}
\begin{itemize}
\item 1.如\hyperref[para:只写失败]{理论部分总结}中所说，知道了怎么正确地识别、分析、处理“失败的策略”，本身就是解决问题的有效方案。仅从内容上来说，本指南已经实现了问题和解决方法的全覆盖。\vspace{-0.5em}
\item 2.如\hyperref[para:想到哪写到哪]{4.6.2小节}中所说，本指南写作时缺乏清晰的规划，内容的编排顺序出了问题，“先叙述解决方案，再叙述问题”的事情频繁出现。这在观感上形成了“罗列问题后没有处理”的印象。\vspace{-0.5em}
\item 3.在本指南中，一部分“解决问题”的讨论篇幅会放在每一章末尾的实操环节中（如本章末尾的“\hyperref[sec:假性成长]{假性成长}”一节）。虽然实操篇章中也会有戛然而止的现象，但总体来说会少一些。\vspace{-0.5em}
\item 4.本指南所能给出的解决方法，和一般的心灵鸡汤没有什么区别。那些方法本身确实有用，但也确实不总是有用。为了准确地阐明其中的细节，我们以下用一个比喻来说明：\vspace{-0.25em}
\end{itemize}
我们将\indicate{行为}（以及行为组/行为模式）比作\indicate{演员}，将\indicate{行为的相互作用机制}（相互触发、并行、竞争和覆盖，以及由它们组成的更复杂的机制）比作\indicate{舞台}。
\begin{explain}
    这个比喻借用自广义相对论。相关的资料中会将“动力学变量”（如“物体的速度”“场的强度”等）称为“演员”，而将“背景时空”称为“舞台”。\footnote{此处主要参考的是梁灿彬先生编写的《微分几何入门与广义相对论》第二版。除了首次提及（上册190页）是借用“惯性质量”与“引力质量”论述“广义相对论为何选择研究时空”外，之后的提及（如中册46页）均用于强调“时空度规的双重角色性”，即“度规既是动力学变量，又是时空背景”。这种强耦合导致了很多困难，比如“解方程难度大幅升高”（上册225页），或是“一些物理学概念失效”（中册46/59/80页）。}
    
    当然，读者理解这一比喻不需要先学会广义相对论\footnote{即使读者熟悉广义相对论，两个比喻也不能直接等同。“行为和相互作用”更贴近于“粒子和场”而不是“粒子和时空”。}。提及它的出处是为了说明：此处使用这个比喻，是想要介绍“如果将演员与舞台混淆，则会大幅度增加分析难度”这一观点，侧重“演员和舞台的对立性”，即“我们将人视为舞台，而将行为模式视为演员”。比喻中不含有“人生就是设定好的剧本和傀儡戏”、“每个人都是逢场作戏”或是“演什么戏自己决定，一定要演得出彩”之类的，关于“演员”和“舞台”的另一些常见指代意象。读者如果更喜欢其它意象，也可自行将其替换为“舞者”和“舞池”等分别指代了“参与者”与“场地”的比喻。

    在接下来的行文中，“舞台”会指代两种不同的东西：一种是“行为的相互作用机制”，后文中也使用\indicate{舞台机制}的称呼；另一种是“具有舞台机制的意识实体”（比如说“一个人”），此时就会有“一个舞台和另一个舞台”之类的区分，后文中也使用\indicate{表演}的称呼。在不引起歧义的情况下，我们会将两种东西都称为“舞台”，读者应能自行区分。
\end{explain}
本指南介绍的理论可以导出一个结论：\indicate{演员可以主宰表演}。
\begin{explain}
对应到术语上，这里指的是“任何一个行为模式的触发范围都可以任意大”。在很多情况下，这会使得“一个人无论面对什么，都只会使用少数几种行为模式来应对”，而这个人的其它行为都无法赢得竞争。这种现象极为常见，在外界看来会体现为“性格”、“能力”等特征。我们在本段中暂时将这样的行为模式称为“主宰行为模式”。读者也可使用“全局行为模式”等自己更顺口的称呼。

主宰了舞台的演员会吸引所有的目光（即“主宰行为模式会得到大量的理论研究”，以至于我们会将舞台当成演员）。我们将这样理论简称为\source{演员理论}。这里的“理论”指的是“对某个行为模式的任何类型的解读”，可以是“一般性的理论”、“私人的生活感悟”等。

由于读者可能从“无法赢得竞争”（即“被覆盖”）这一表述中读出歧义，此处需要特别说明一下“潜意识”。对应“潜意识”这一概念的，并不是“被主宰行为模式覆盖掉了的其它行为模式”，而是“主宰行为模式的（不自知的）成因”。只有后者才能起到“觉知了潜意识后，放下了心结，得到了成长”的效果。由于“主宰行为模式经常是由这些被覆盖的行为模式演化而来的”，所以这二者经常是同一件事，但是它们不总是相同。

另一方面，被覆盖的东西不一定是不自知的，成因也不一定是不自知的，它们和潜意识是三件不同的事。这也是本指南极少使用“潜意识”一词（并且提及时基本都持负面态度）的原因之一：它的定义较为模糊（并且随着广泛的社会传播，歧义很重），并且不是本指南理论建构的必要组件。
\end{explain}
演员实际上还在按照舞台机制行事，于是，\indicate{观察演员时，也可以得到一些舞台的性质}。
\begin{explain}
如果读者喜欢的话，也可以使用一些更文学化的表达，比如说“我们每个人的身上都有外部世界的完整投影”、“集体潜意识塑造了每个人的生活”之类。

“透过演员来分析运用舞台机制”的理论随处可见：一些是深奥的理论，一些是脍炙人口的俗语（比如“一千个读者就有一千个哈姆雷特”），一些是充满智慧的人生格言（比如“时间就像海绵里的水，挤一挤总会有的”），一些是简单的生活小妙招（比如“看到对方难受了就赶紧转移话题”）。

确实有一些演员理论（其中一些在运用了大量的隐喻、象征和指代后）能够完整地刻画舞台机制。很多个不同的心理学理论都能够解释所有的现象。

一些很常见的，会被用来解释行为和动机的东西，其实都应被视作演员：生理反射、爱、恨、情绪、状态、冲动、自我防御、潜意识、自恋、执念、俄狄浦斯、利益、等价交换、道德、价值观、理想......它们可以用来刻画舞台机制，但还是在\indicate{试图让单一演员操控整个舞台}。%它们只是“一种解释”，而不是“唯一的解释”。

这会带来两方面缺点：一方面是“舞台机制通过理论的间接描述体现出来”，这会增加理解和运用的成本；另一方面是“难以实际操作”，象征带来了大量的模糊解读空间和歧义。比如，如果在将理论按照字面理解，那么就只能依靠（自己或者某个其它人身上的）一个演员来操控整个舞台。虽然它在别的舞台上是主宰，但这个舞台上可能非常边缘，甚至完全不存在。如果想强行将其推向台前，则容易产生更不利的结果（详细讨论见\hyperref[sec:多行为的形成]{5.4.4小节}）。
\end{explain}
同时，\indicate{观察一个演员时，会忽视其它演员}。
\begin{explain}
对应到术语上，这里指的是“研究一个行为模式时，会无法稳定地得到其它行为模式的性质，大量的性质不会在研究过程中出现”。对“一个行为模式”来说，它的“构成”（即“都包括哪些行为和行为链”）和“成因”（即“是经历了什么才成这样的”）是比较好研究的两方面。

对于每一种常见的行为模式，我们都能提取出一个模板，比如“教科书里的知识”、“一类生理反应”、“某些刻板印象”等。但具体到每个人身上，行为模式的构成总是会和模板有一些差别。这些差别是在形成过程中，混进来的一些“其实不属于这一行为模式”的东西，因人而异，毫无规律可言。

如果我们不将其视为“舞台的整体性质”，而是想要从“演员的性质”出发，推导出“这些差别（意外情况）存在的合理性”，那么如果不将演员庸俗化变成“万物皆有可能”，就容易发现“意外情况好像无穷无尽，怎么都描述不完，总有漏网之鱼”。这容易导出“语言（认知）是有极限的，实在界总是有不可言说之物”一类的结论。但是，很多情况下，只要我们换个视角（此处指“换一个演员看”，即“观察另一个行为模式”），有时甚至只是“多起了一个名字”，就能清晰地完全描述之前视角下论述不清的东西。
\end{explain}
本指南持有的观点是：\indicate{舞台是存在的，并且舞台机制主宰了表演}。
\begin{explain}
对应到术语上，这里指的是\indicate{意识现象客观存在，遵循其客观规律；其客观规律真实存在}。“意识现象”与\hyperref[sec:人的意识演化基本模型]{之前的定义}一致\footnote{严格来说，此处的“意识现象”不仅可以指“人”，而也可以指代任何“产生了竞争与覆盖机制”的实体。但如\hyperref[sec:澄清与叠甲]{3.6.2小节}中的讨论所说，以本指南的篇幅，显然无法涵盖对社会、人工智能等领域的讨论。读者也不应忽略这些领域的特点，将本指南的分析框架生搬硬套到其它领域去。}，也是“行为及其相互作用机制”，也即“演员和舞台”。

此处的“客观存在”，指的不是“它有对应的物质实体”，而是“它是一种稳定可复现的现象”。“意识现象及其原理”是纯粹的“软件理论”，“行为”本身就来源于神经现象的复杂涌现，不需要（并且在很多情况下也不可能）将其还原至物质实体的层面；“行为”本身就会参与舞台中复杂的演化，不应将所有行为都还原为外部输入。

概括地来说，演员（即“行为”）所体现的是“因果性什么时候有效”（因为行为本身就是对因果性的刻画：刺激是因，而行动是果），而舞台机制所体现的是“（某一特定的）因果性什么时候无效”（被覆盖时，因果链就被打断）。二者共同组成了对因果性的完整刻画。
\end{explain}
\indicate{本指南所介绍的理论，是关于舞台的理论}。
\begin{explain}
我们将其简称为\indicate{舞台理论}。更进一步地说，本指南认为舞台机制（即“相互触发、并行、竞争、覆盖”）是固定不变的；而人（表演）的变化，完全由“演员的变化”所体现。

这种视角会带来一些明确的缺点，其中最明显的大概是“这套理论完全无法预言‘舞台上会有哪些演员’”。演员的全部信息，完全需要从外界获得；而这些演员的独角戏（即“行为模式的触发和其内部的演化”），以及演员之间的对手戏（即“行为模式之间的竞争和演化”）会如何进展，也经常需要从外界获得（比如某门学科的研究课题，或某个人的具体表现）。舞台理论能做到的只有“已知舞台上有哪些演员时，分析这些演员会怎么演出，又会有哪些新演员”，即\indicate{在已知所有行为和它们的竞争关系后，就可以推算它们的演化}。

\trained 或许会从上面这些论述中发现一些哲学观点。但本指南无意将这一理论上升到哲学层次，也推荐读者优先从本指南给出的定义和分析出发，理解这一分析框架。

如果读者认为本指南的理论与某个另外的理论相似，那么无论是认同这一理论还是排斥这一理论，都将影响对本指南所介绍理论的准确理解。虽然这是一套舞台理论，但承载舞台理论的认知仍然是演员，仍然会和其它理论展开竞争和覆盖。读者可能可以注意到，本指南在行文中极少使用专属于哲学领域的概念，或是一些广泛讨论和传播（从而也有很多歧义）的概念。这就是为了防止自身的理论受到这些词语干扰。对于熟悉这些概念的读者，不使用它们会增大理解难度。本指南对此深表歉意。
\end{explain}
舞台理论所关注的事情，首先是“没有遗漏地\indicate{描述}所有的演员，以及这些演员的互动”，而不是“分析这些演员的互动”。
\begin{explain}
\indicate{对所有演员的分析，就是对舞台的完整分析}。无遗漏性会使舞台理论体现出一定的分析能力，但这仅代表“它可以运用从已知现象中总结出的规律”。分析能力不属于舞台理论本身，只是对演员理论的借用。若现象不足以总结出规律（如“案例太少”、“没有代表性”、“只有无关信息”等），则舞台理论毫无分析能力。

无遗漏地描述所有的演员，即“了解一个人的所有行为以及它们的竞争关系”，这一条件听起来完全不可能实现。但在实际操作中，通常只需要关注少数行为，即“某一行为模式的\note{有关现实}”即可。其它的行为如果不能打断（与其并行，或者被其覆盖，均可）它，那么就相当于不存在。这样，我们就得以在一个小得多的舞台上\footnote{\label{footnote:局部}这一操作可以称为“舞台的局部化”。“局部化”是数学领域的借词，这一思想在很多方面（比如说“函数芽”等）都有所应用。与此处最像的例子是，在研究“某个环的某一极大理想”时，构造一个“与该环性质相似，但仅包括唯一的极大理想”的环（这个环被称为“局部环”）。读者理解这一借用不需要先学会深入的环论和代数数论，仅需理解“这是一种突出一个演员的方法”即可。\hyperref[sec:局部有效理论]{5.2.3小节}中，“局部有效理论”的“局部”也因此而指“研究单一演员的理论”。但“局部”这一概念不是很精确，同时也不是理论的必要组成部分，所以不正式引入。}，研究单一演员的演化。
\end{explain}
\indicate{在关注“一个演员与其有关现实”时，舞台理论与演员理论不可区分}\footnote{\trained 可能会更习惯“在局部，舞台理论退化成演员理论”的表达，但是“退化”一词含有“理论之间的价值判断”，所以此处不使用这一表达。这些读者应该也知道“理论A变为理论B，说明理论A没有引入新的本质困难”这一结论。对应到此处，就是“舞台理论相对演员理论的困难，全部集中于‘如何确定有关现实’这一点上”。}。
\begin{explain}
此处的“不可区分”指的是“经过完全相同的分析过程，得到完全相同的结论”。

这使得我们得以借用关于演员理论的所有内容，比如说“所有类型的成功经验”。如果一个行为模式是有效的，那么就可以尝试将其装载到自己身上。也是因此，本指南能给出的解决方法，和其它资料中的解决方法并没有什么区别。而且，本指南由于的篇幅和话题限制，不可能对所有问题展开详尽的讨论，其它资料中的解决方法会更加详细可操作。

舞台理论和演员理论最大的不同在于，在不确定“演员和有关现实”的情况下，我们不知道要将舞台理论变为哪个演员理论。从舞台理论中可以得到“任何方法都不保证有效”、“任何行为（指一对‘刺激-行动’）都有可能出现”等结论，这说明在尝试达成任何一个目标时，都不能事先选定一个演员理论。

这一思路可以导出很多有效的操作思路，比如：先根据目前的信息制定计划，如果执行不下去，遇到了新的阻碍，再针对性地调整计划，处理阻碍。这实际上就是“切换了一种演员理论”。
\end{explain}

\section{实操：假性成长\label{sec:假性成长}}
\hfill\begin{minipage}{0.6\textwidth}
\fontsize{8pt}{12pt}\selectfont\fontsize{8pt}{12pt}

\raggedright 或许廉价的尊严正在悄悄瓦解，或许该杀死我丑陋的自卑，再变得虚伪。\footnote{\bilibili{BV1m44y1i7ri}。}

\raggedleft 鸽稀拉《狼狈》

\raggedright 所谓成长，不过是将心血滴成冰冷的离别。\footnote{\bilibili{44463232}\another\netease{1349566322}。}

\raggedleft 时之歌project\&苍白《花虹》

\raggedright 大人说笑一笑解千愁，可是为什么，为什么，心中有火？\\
是不是我还太小，境界还不够？\footnote{\bilibili{av6107357}\another\netease{1377097873}。\\}

\raggedleft 绛舞乱丸《寿星街小结巴》
\end{minipage}

在绕了好大一圈之后，我们终于可以来讨论\hyperref[para:假性成长]{第三章一开头}提到的问题了：为什么明明听过很多道理，却依然过不好这一生？
\begin{explain}
现在，我们已经可以较为轻松地描述这一现象了：那些道理无法赢得\paranote{可能是“完全没有参与”，也可能是“参与了但被覆盖”。}相关的竞争\paranote{具体竞争可能有“克服困难”、“改正坏习惯”、“学会新能力”等等。}，无法形成一个有效\paranote{此处的“有效”根据目标而定，目标可以是“考上好学校”、“开心”、“赚大钱”等等。}的行为模式。\vspace{0.5em}
\insertparanotes
\end{explain}
为了更为准确地描述和处理这一类现象，我们引入如下的概念：

在一个环境中，我们将\indicate{两个行为模式的有关现实相同}的现象称为“这两个行为模式在这一环境下\indicate{同步}”。对应地，如果\indicate{两个行为模式的有关现实不相同}，我们就称“这两个行为模式在这一环境下\indicate{不同步}”。
\begin{explain}
与之前一样，这两个行为模式可以在同一个人身上，也可以分属两个人。

两个行为模式间的“同步”性质要比“不可区分”性质更弱（即“不可区分一定同步，同步不一定不可区分”），这主要是因为两点：
\begin{itemize}
\item 同步仅限制了“触发”和“打断”相同，没有对行为模式内的竞争加限制。不过，我们可以在有必要的情况下加强定义，规定“行为模式内部的竞争结果也需要对应”。这不会影响接下来的讨论。
\item 同步不要求“行为之间不可区分”，行为可以有很明显的差别（比如说“一个是实际行动，而另一个是脑内的模拟推断”）。这是这两个概念更根本的不同，它允许我们刻画“一种行为模式和对它的认识之间的联系”之类，有紧密联系但是组成不同的行为模式。

同时，同步也不要求“它们会被相同的刺激触发/打断”，而只是要求“它们会同时触发”。即：“两个行为模式可能因为同一事物的不同方面特点，而同时分别触发”一类的现象也算作行为模式同步。但是，“因为同一环境内的不同事物而触发”则不算同步。其判断标准为“两种不同的刺激是否可以分离，是否可以独立出现”，这可以避免我们对特征无限细分。
\end{itemize}
“同步”是一个与环境相关的概念，这么定义会方便我们描述现象与简化讨论。两个行为模式可能在大部分环境中都同步（比如说“完全不触发”，或者是“在简单环境中都能解决所有问题”等），只在少部分环境中不同步（比如“死板地照猫画虎，在需要灵活变通的时候就没办法了”等）。我们只关注这些不同步的少部分环境。
\end{explain}
如果在行为模式A的触发环境中，行为模式B不与行为模式A同步，那么行为模式B就是行为模式A的无关现实。
\begin{explain}
此处为了使得结论更为简短，略去了很多关键细节。请读者将“环境”具体到“每一个会触发行为模式A的现象”上，分别判断是否同步。单一现象下，“不同步”与“无关现实”确实等价。之后的行文中也会按照这种处理方式（即“具体到单一现象”）使用“无关现实”这一概念。

因为我们增添了“A的触发环境”这一条件，所以“B和A不同步”仅有两种可能：
\begin{itemize}
\item A触发的时候B不触发。这种情况下B自然不参与任何竞争，是A的无关现实。于是，这不是我们接下来讨论的重点。
\item A触发的时候B也触发，但成因不同。此时，B不仅没办法解决A的问题，还会产生新的问题，从而更糟。
\end{itemize}
“B触发而A不触发”的情况不在此处的讨论范围内。这种情况当然也存在，比如说“在A不触发（比如说‘空闲时间’、‘提前准备’）时，分析‘应该怎么应对A’，并得到一些结论”。

我们暂且将“A不触发”称为“场外”，而将“A触发”对应地称为“当场”；而这种情况称为“B（在A的）场外触发”。B有可能同时当场触发，也可能“只在场外触发，但‘分析得到的那些应对方式’会当场触发”。但总之，我们只研究当场触发时的情况。
\end{explain}
到此为止，我们已经能够完全描述这一现象。我们以一个例子作为引入：
\divider
有很多学生每天学习很长时间，但成绩却没有什么起色。有些分析会将其称为“\indicate{假努力}”/“\indicate{无效努力}”。
\begin{explain}
本节的标题“假性成长”就是从这个概念中引申出来的称呼。

“无效努力”一词是针对结果的，只要成绩没有起色，那么用这个词就没什么问题。而“假努力”则更多地指“表演性质的努力”，会指代一些具体的行为。我们以下将分析重点放在这里。
\end{explain}
关于无效努力，一种常见的批评是：“只会刷题/记笔记，但不会整理反思，所以学了就忘”。根据实际情况的不同，这一批评能起到的效果也很不同：
\begin{explain}
为简单起见，我们此处仅分析那些“有明显套路，可以被很多学习好的人完全分析透”的题\paranote{比如一些基础模型和二级结论。}。同时，A有“包含这些套路”的资料，比如说“自己的思路和过程”\paranote{一些学生可以仅从“自己做错了题”这一个信息中，就检查出思路和过程的所有问题。}、“标准答案”、“老师专门讲的题（以及抄下来的笔记）”等内容。

在下面的例子中，我们将“无效努力的人”记作A，将“给A提供教学（或其它方面帮助/评价）的人”记为B（和之前一样，此时我们也不要求A和B属于不同的人）。我们会按照编号将下面这些分别称为“情况1/2/3/...”。
\begin{itemize}
\item 1.确实有一些A完全只是没意识到要这么做。这样的A在听到了B的这句话\paranote{不包含任何其它内容，就真的只是“你不会整理反思”这一句话。}后，就能立刻明白要怎么做，并且立刻做到位。A在“怎么做”上完全不存在任何的\note{信息缺失}，于是可以立刻培养出对应的行为。

这个过程符合简单的“知错就改”\paranote{“错”指“不反思”，而“改”指“会反思”。此处说的不是“‘错’指‘题做错了’，而‘改’指‘看标准答案’”。}\note{演员理论}。我们完全不需要退回到舞台上，只需要关注这一演员，就可以解决问题。

\item 2.一些人除了没意识到要这么做以外，还不知道要怎么做。在一些情况下，只需要B简单地教学\paranote{比如“B告诉A：你要主动去想‘标准答案里面都用到了什么知识’、‘你自己又是怎么没想出来’等一系列的问题”，或是“教学资料中包含对思路的分析，A只需要跟上就行”。}，A就能完全学会。这样的教学明确给出了关注点，这些关注点补足了信息缺失。

大多数情况下，B只要自己知道“怎么整理反思”，那么就会顺口说出这些话。虽然这里已经涉及到了一点点分类讨论的内容\paranote{不是所有人都需要这些唠叨，比如，情况1的A就不需要。}，我们也不会将其视为“回退到了舞台上”，而仍然会认为“这是‘反思’这一技能的必要知识点，是‘知错就改’中的必要步骤”。
\end{itemize}
\insertparanotes
\end{explain}
但其它的情况就没这么简单了：不仅是A有信息缺失，B也开始有信息缺失了。
\begin{explain}
\begin{itemize}
\item 3.一些A即使听到了“标准答案里面都用到了什么知识”，也看不出来标准答案用了什么知识；或是“表面上能看出来用了什么知识，但是总结不出来思路”。在这种情况下，无论B强调了多少遍“要整理”，都没有用。

不过此时，我们可以轻松地后退一步，换一个演员理论。会教\paranote{比如有常年的教学经验，或是看了很多相关的讨论、分析、经验分享。}的B可以轻松地看出来，这是因为A有更基础的知识漏洞。“（即使被细致地教了，仍然）看不清思路”、“无法迁移”等问题，是因为“从一开始就没有理解这些思路”。此时需要全面地复查“A都有哪些知识漏洞”并修补，而这可能会耗费非常大量的时间和精力\paranote{比如说“已经高三了，但必须从初二开始补”之类的情况。}。

当然，此时也会有“只是一节课（一个月）没听，所以不知道某个技巧”之类的简单情况，这可以用“查漏补缺”这一演员理论来解决。
\insertparanotes
\end{itemize}
\end{explain}
而当知识漏洞更大时，情况远远不止“会花更多时间和精力”这么简单。
\begin{explain}
为简便起见，我们以下假设“B有无限耐心，可以把每个问题都讲懂；A也无限专注，能够记住B之前讲的所有东西”。
我们此处不讨论“消耗精力”的方面，只关注一类无效的处理方式：
\begin{itemize}
    \item 4.让A自己补。它有多种不同的变体：
\begin{itemize}
\item 更加详细地讲题：从一道A学不懂的题开始，B详细地讲解每一个不理解的地方。如果讲出来的东西A仍然听不懂，那么B就继续讲更基础的内容。此时，结果也很容易变成“A除了这道题以外还是什么都不会”。

这是因为，B仅能针对“A在这道题上的疑问”展开讲解\paranote{此处我们不讨论“B开始给A从更低年级的知识开始讲”的方案，这一方案并入下一种情况讨论。}。这只能包含“这道题所需的思路”，不包含“这道题用不上（但是其它题能用上）的思路”。这些思路必须要靠补基础才能练出来。A不知道“什么时候该用什么方法”，这是不可忽略的信息缺失。B没有意识到，也没有处理这一信息缺失，所以自己的思路必然不可能和A的思路同步。

\item B直接给A讲基础知识：我们此处假设B的教学资源无限丰富。此时需要面对的主要问题是：A容易被其它东西分心，比如“自己不会做的题\paranote{尤其是，如果B是从一道“A不会做的题”引申过来的（即上一种情况中并到此处讨论的内容），那么A就很容易频繁思考“讲的这些到底怎么才能运用到刚才那道题上”之类的问题。A如果没有意识到“这些知识要补几天甚至几个月”，并且据此来（在学完之前）放弃“运用到刚才那道题上”的尝试，那么就会因分心而漏掉很多知识。}”、“其它的学习任务”等。这些东西对A来说，明显会“更重要”、“更应该做”。

A此时不一定能够意识到“补基础”的必要性和重要性。这主要有两方面原因：一方面是“A可能对基础的理解有误”\paranote{比如将“巩固基础”完全理解成了“多刷题”。A确实也有可能靠“背诵某些固定套路”的方式，在低年级成绩还不错。但A从未学到思路，也不认为B讲的思路是什么值得留意的东西（参见情况3）。}，从而产生了信息缺失，无法判断自己哪里有不足，又怎么算补上了；一方面是“A无法形成B的思路”，无法运用\paranote{此处的“运用”指的是：A可以轻易地承认这一结论，但是难以形成“确信补基础要比其它事情更重要，从而把其它事情的时间腾出来，优先补基础”的完整思路和决策。}“补基础真的很重要”这一事实。
\insertparanotes
\end{itemize}
此时我们会发现，要讨论清楚这种情况，就必须引入更多的有关现实，必须讨论舞台机制，而不能只使用“查漏补缺”的演员理论。每一个局部都是很合理的，但是当它们拼到一块，就无法相容。\end{itemize}
\end{explain}
而B如果无法精确地得出“这是一个必须用舞台机制分析的问题，不能只盯着学习看”这一结论，就会得出另一个结论：“A的学习态度有问题”。
\begin{explain}
确实如此，这是一句正确的废话。“假努力”本身也是学习态度有问题。但是，形成了这个念头，就会产生另一种处理方式：
\begin{itemize}
    \item 5.让A自己改。这几乎是不可能实现的。
    
    很多人会把情况1和情况2当成是“A学习态度良好，一说就改”的正面典型，而把情况3\paranote{仅指那里提到的“可以靠查漏补缺解决”的情况，在外部的表现可能是“一次没考好，但是后面就恢复了”，或者是“之前成绩差，花了一个学期把成绩提上来了”之类的。}当成是“A亡羊补牢，为时不晚”的正面典型。但这样的做法不可迁移，因为客观条件不同。

    此外，我们确实可以找到不少“因为态度改变了，所以能学习了”的例子。这样的例子通常是因为“接触了一些新东西\paranote{比如说“一位好老师”、“惨痛的社会教训”等等。}，得到了一些新的感悟”。这又是一些“A浪子回头，迷途知返”的正面典型。这种情况下，A会对“自己为什么能重新开始学习”有清晰的认识，并且讲\paranote{此处不区分A是“会主动讲”还是“别人问到了才说”。此处仅关注“大家都得到了‘态度变了就能学好习’的认识”这一点。}给其它人听。

    读者此时可以发现，上面的这些内容，其实已经不能用单一的演员理论涵盖了。虽然我们仍然可以将其全部称为“态度转变”，但它们所面对的竞争和产生的覆盖\paranote{比如，“自己态度改变的原因”很多时候，起到的真正作用是“覆盖了之前的行为模式”，而不是“提供了正确的学习方法”。如果另一个人没有同样的行为模式要覆盖，那么这个原因就没有用。}都完全不同。除了“作为成功学案例激励\paranote{激励确实只能起到激励效果。有少部分人会因为这些案例而认真学习，但是大部分人没有什么改变。这些人也会觉得“成功的人很厉害，我应该学”，但是没什么用。}很多人”之外，这样的案例不会起到什么作用，不提供任何可迁移的经验。

    更糟糕的是，在A想自己改的时候，还会有很多其它行为模式来和它竞争。比如，如果A同时排斥学习\paranote{这在“学不进去”的人的一种常见状态，也是“学不进去”的主要原因之一。排斥的方法有可能是“看到题目就烦”、“大脑一片空白”、“立刻联想到玩”、“焦虑然后自我攻击”等等，因人而异。}，那么A“将自己强压着学习”的行为，本身就会打断一部分收集和整理信息的行为\paranote{比如说“记忆”、“读题”、“串联思路”等等。}。这会形成恶性循环，导致越努力越学不进去。
\end{itemize}
\insertparanotes
情况5的高失败率，源于“A需要将某一个行为模式整体更换为另一个行为模式”。而无论是“态度有问题”这个判断，还是那些零散而无规律的成功经验，都既不能提供“原来的行为模式哪里不行”的信息，也不能提供“新的行为模式应该是什么样”的信息，与真实情况处处不同步，完全没有办法提供任何参考。信息缺失时只能向后竞争，不可能稳定成功。
\end{explain}
\divider
在上面的例子中，我们不可避免地分了很多类，讨论得相当琐碎。这一特征在本指南的一些其它地方（比如\hyperref[sec:情绪、感受与其分析与处理方法]{3.7节}、\hyperref[sec:对竞争过程的解读]{4.4.3小节}、\hyperref[sec:指挥现象的一个实例解读]{4.5.4小节}、\hyperref[sec:日常沟通与社交]{4.7节}、\hyperref[sec:局部有效理论]{5.2.3小节}）也有体现。这能够涵盖所有主要的情况。

但即使这样，应该也有很多读者会觉得“此处的讨论仍然不够贴合自身的实际经历和体验”。一方面这是因为篇幅限制，想讨论的话会有无限的细节可以展开（比如每种排斥学习的态度都能展开讨论，而态度有无限多种）；另一方面，如果讨论得过于详细，就会包含大量“原因相同，结果相反”的内容，反而会出现“完全与读者的体验背道而驰”的现象。

本指南此处所采用的讨论方式是“尽量全面地讨论某一个行为模式以及它的所有有关现实”（也就是“将舞台局部化”），但有关现实仅是提及，而不做进一步讨论。不同的舞台的有关现实不同，进一步的讨论已经不应该由同一个演员承担，而是应该转换视角，独立考察另一个演员。
\divider
我们从以上的分析中，抽象出以下定义：

对于一个目标，如果\indicate{某个“因这一目标形成的意识现象”不能实现这一目标}，那么我们就将这一意识现象称为这一目标的\indicate{假性成长}。如果意识现象有所指代，那么我们就将其指代的内容称为这一目标的\indicate{假性经验}。
\begin{explain}
此处的“意识现象”可以指代“某个认识”、“某个行为模式（能力）”等很多现象。“成长”一词指的是“有新的意识现象形成”，这在一些时候会给人带来成长感\paranote{不是所有新的意识现象都会给人成长感，有些意识现象（比如“习得性无助”、“记忆封闭”等等）带来的是挫败感，有些意识现象（比如“坚持努力”、“宽慰自己”等等）会转移注意力。但它们也符合“假性成长”的定义，也会被算作假性成长。}。

任何客观现象都可以成为“目标”。它可以指某个外在的追求，也可以指某个意识现象\paranote{如果是某个难度较高，较为长期和遥远的外在追求，涉及到了“策略制定、行为形成与能力培养”，那么也可将目标视为意识现象。}。如果它指某个意识现象，那么此处定义中，要求“指定该意识现象”。这在一些情况下很好描述\paranote{比如说“感到开心”就只涉及某一特定行为/行为模式。}，一些情况下需要费一些力气\paranote{比如“自我实现”经常要求“通过某一种特定的行为模式获得一个外在的成果”。}，还有一些情况很容易描述错\paranote{比如说“不痛苦”在很多时候无法用“感到开心”、“看开点”之类的内容替代，因为它们经常无法完全打断“痛苦”这一复杂的行为模式（或许可以打断一些行为，但这无济于事），而只是和痛苦并行。}。
\insertparanotes
\end{explain}
当目标指代意识现象时，此处的定义不要求“目标和假性成长属于同一个人”。
\begin{explain}
这一处理主要是为了方便读者理解。当它们属于不同的人时，虽然“假性成长”这一称呼会有些奇怪，但是更容易展示现象：

我们延续之前的记号，即“目标在A身上，而假性成长在B身上”。此时的情况类似于“B作为过来人，给A讲述大道理、正确做法和人生经验”。此时，“假性成长”指的是B“觉得A应该这么做/那么做不好”一类的认知和反应，和“对A的评价、指导和教育”；而“假性经验”指的是B“亲身体会”、“从别处听来的道理”、“觉得这么做好”的认知等内容。

行为模式是属于不同的人还是属于相同的人，其实对分析没有区别。我们仅需注意“两个行为模式是否互为有关现实，是否同步，是否能成功地竞争，有效地打断和覆盖”。如果目标和假性成长属于同一个人，那么上述的现象依旧会发生。A仍然有可能一边觉得“自己应该这么做”，一边又完全应用不到实际情况中去。
\end{explain}
假性成长不能实现目标，是因为它不能有效地参与到竞争过程中，不能消除原先的行为模式，形成新的目标行为模式。
\begin{explain}
这不是一句废话。它用于说明“这是一个只能用舞台理论分析的问题，无法使用单独的演员理论\paranote{使用单独的演员理论就会有无穷无尽的反例。}”，具体需要参与什么竞争过程只能依靠当场调研，每个舞台都不一样。越是长期整体的目标，就会接触到越多的有关现实，就需要越全面整体地考虑舞台上的全部演化。

不是所有的目标实现都需要打断原先的行为模式，比如“在空闲时间接触一个新领域”就经常没什么需要打断的\paranote{但我们也会看到“一有空就玩”、“一有空就内耗”、“一有空就做家务”等多种反例，此时则需要打断。此处的“空闲”指的是“没有（主宰）行为模式触发”，但如果空闲本身也可能触发另外的行为模式。}；不是所有的目标实现都需要形成新的目标行为模式，比如“改正一个坏习惯”就不一定需要用新行为模式来对冲\paranote{改坏习惯可能只需要克制，可能需要形成几个零散的行为，可能需要从根上彻底重塑生活节奏。不同的情况需要做的事情不同。}。
\insertparanotes
\end{explain}
假性成长确实也形成了某个行为模式（或者是它自身，或者是它和假性经验），但因为无法参与竞争，形成的行为模式完全是随机的，比如：
\begin{explain}
\begin{itemize}
\item \indicate{被污染}：如果想要消除原先的行为模式，但假性成长无法在竞争中取胜，那么因为形成了关联，原行为模式就会外溢到“想要改变的尝试”上\paranote{比如，定了闹钟但还是会睡回去，久而久之闹钟就叫不醒了。}。这会使得下一次的改变更困难\paranote{原先“可以并行”的刺激（和对应行为）也被覆盖了，少了一个可用的选项。}。
\item \indicate{变心}：假性成长可以和原行为模式并行，有可能在一些情况下占据优势。随着时间推移，原行为模式因为其它的原因而消失，只有假性成长保留了下来。假性成长无法形成新的行为模式，变成了盲目的执念\paranote{如果执念会造成比较严重的问题，那么就需要考虑消除它。此时，“通过潜意识，来觉知自己当初为什么会这么想”的方法可以回想起原来的目标。但回忆也只能参与向后竞争，不一定能消除这一执念。}。
\end{itemize}
\insertparanotes
\end{explain}
读者应该也可以注意到，所有关于假性成长的讨论中，都有“一个目标”作为前提。这是不可缺失的条件。如果完全没有目标，那就不可能提及“成长”\paranote{此处不应将“目标”功利化（比如“我不要什么目标，我只要开心”）。“目标”的定义为“任何客观现象”（包括意识现象），可以包含所有想得到的讨论内容。}。目标由舞台的演化提供\paranote{比如“儿时的兴趣”、“生活所迫”、“假性成长”都可以提供目标。}。本指南的分析框架中不包含任何“提供目标”的部分\paranote{但包含“将一个目标转化为另一个目标”的部分，比如说\hyperref[sec:能力的形成]{行为形成}的篇幅。}，只包含“目标的后续处理”。
\insertparanotes